% v4.tex -- Pedagogical report on the sBootstrap framework
% Clean rewrite from physrevD.tex (verified core), v4_patched.tex
% (pedagogical expansions), v3.tex (worked examples), and session 28 notes.
% 12pt article class, ~36 pages. All derivations shown.
\documentclass[12pt,a4paper]{article}

\usepackage{amsmath,amssymb,amsfonts,amsthm}
\usepackage{graphicx}
\usepackage{hyperref}
\usepackage{booktabs}
\usepackage{multirow}
\usepackage{geometry}
\geometry{margin=2.5cm}

\newtheorem{theorem}{Theorem}
\newtheorem{proposition}[theorem]{Proposition}
\newtheorem{corollary}[theorem]{Corollary}
\newtheorem{lemma}[theorem]{Lemma}

\newcommand{\SU}{\mathrm{SU}}
\newcommand{\UU}{\mathrm{U}}
\newcommand{\SO}{\mathrm{SO}}
\newcommand{\Sp}{\mathrm{Sp}}
\newcommand{\Nf}{N_{\!f}}
\newcommand{\Nc}{N_{\!c}}
\newcommand{\Tr}{\mathrm{Tr}}
\newcommand{\adj}{\mathrm{adj}}
\newcommand{\diag}{\mathrm{diag}}
\newcommand{\keV}{\,\mathrm{keV}}
\newcommand{\MeV}{\,\mathrm{MeV}}
\newcommand{\GeV}{\,\mathrm{GeV}}
\newcommand{\TeV}{\,\mathrm{TeV}}
\newcommand{\ie}{i.e.\@}
\newcommand{\eg}{e.g.\@}
\newcommand{\zch}{\mathfrak{z}}

\title{From mesinos and quarkinos to the Standard Model:\\
  the sBootstrap as a structural EFT from Seiberg confinement,\\
  charge--mass duality, and generation uniqueness}

\author{A.\ Rivero\thanks{Corresponding author. \texttt{rivero@unizar.es}}}
\date{Draft v4 --- \today}

\begin{document}
\maketitle

\begin{abstract}
We identify Standard Model fermions with composites of a
confining $\SU(3)_F$ family gauge theory: mesinos
($\psi_M$ of $M^i{}_j = Q^i\tilde{Q}_j$) are leptons;
adjugate composites provide quark partners via partial
compositeness.
The charge-squared mass law $m_i = \mu\,\zch_i^2$ is the
unique monomial for which $K = \sum\zch_i^2/(\sum|\zch_i|)^2$
is independent of the Cartan direction (Theorem~1);
the Seiberg seesaw then produces both a direct branch
($K = 2/3$ for leptons, $0.001\%$) and an inverse branch
($K^{-1} = 2/3$ for down quarks, $0.22\%$) from a single UV
condition.
The confined spectrum automatically produces a pair
$(M_i, \adj(M)_i)$ per family with the mathematical structure
of an $\mathcal{N}{=}2$ hypermultiplet---emergent from
confinement plus the adjugate identity, not from UV symmetry.
An O'Raifeartaigh seed at $gv/m = \sqrt{3}$ breaks SUSY with
$K = 2/3$ exactly; the bloom is driven by $\cos(3\delta)$
from $\mathbb{Z}_{6}$ instantons.
The quark chain
$(t,b,c)$--$(b,c,-s)$--$(c,s,0)$--$(s,0,m_u{+}m_d)$
has $K$ within $1.2\%$ in every link (joint $p \sim 1/21\,000$),
and is compressed by three algebraic relations plus a
completion formula that predicts $m_t = 168.5\GeV$ ($2.3\%$).
Five independent criteria select $N = 3$ generations, including
a Casimir--Koide identity from Kostant's principal
$\mathfrak{sl}(2)$ embedding.
The $Q_2$ operator connecting the two branches is the
spin-2 component of $\mathrm{End}(V)$ under the principal
$\mathfrak{sl}(2)$, satisfying a cubic algebra
$R^3 = \Tr(D^2)\,R + 2\Lambda^6\,I$.
The principal open problem is the non-perturbative K\"ahler
metric, which controls the mapping from holomorphic to physical
masses.
\end{abstract}


%======================================================================
\section{Introduction}
\label{sec:intro}
%======================================================================

The Standard Model contains 13~free parameters in the flavour
sector---six quark masses, three charged-lepton masses, three CKM
angles, and a CP phase---with no known relations among them.
While the gauge sector is tightly constrained by the symmetry
group $\SU(3)_c \times \SU(2)_L \times \UU(1)_Y$, the Yukawa
couplings that determine fermion masses are arbitrary complex
numbers.
The absence of flavour structure is one of the outstanding
puzzles in particle physics, closely related to the hierarchy
problem (why $m_e/m_t \sim 3 \times 10^{-6}$?) and to the
origin of CP violation in the quark sector.

This paper develops a programme that addresses the flavour
problem not by postulating textures or discrete symmetries on the
Yukawa matrices, but by \emph{identifying the observed fermions
with composite fermions} of a confining gauge theory.
The key idea is that the mass hierarchies and mixing patterns are
not free parameters but consequences of the non-perturbative
dynamics of confinement: the Seiberg seesaw, the adjugate map,
and the charge algebra of $\UU(3)_F$ together constrain the
mass spectrum.

\paragraph{Historical roots.}
The idea that hadrons are composites of more fundamental quarks
was developed through the ``bootstrap'' programme of
Chew~\cite{Chew:1962} (Rev.\ Mod.\ Phys.\ \textbf{34}, p.~394),
which treated all hadrons democratically as bound states of one
another.
In 1966, Miyazawa~\cite{Miyazawa:1966} (Prog.\ Theor.\ Phys.\
\textbf{36}, eq.~1) observed that the spectrum of hadrons
exhibits a \emph{supersymmetric} pattern, with mesons
($q\bar{q}$, bosons) and baryons ($qqq$, fermions) forming
approximate supermultiplets.
Catto and G\"ursey~\cite{CattoGursey:1985} (Nuovo Cim.\ A
\textbf{86}, \S2) sharpened this: the diquark $[qq]$ (in the
antisymmetric colour $\overline{\mathbf{3}}$) transforms as
the SUSY partner of the antiquark $\bar{q}$, so the baryon
$q[qq]$ is literally the fermionic partner of the meson
$q\bar{q}$.

\paragraph{Luty--Mohapatra models.}
Luty and Mohapatra~\cite{LutyMohapatra:1997} (Phys.\ Lett.\ B
\textbf{396}, p.~161) constructed explicit SUSY composite models
where all quarks and leptons emerge from confining $\SU(N)$
dynamics, with quarks as mesons and leptons as baryons.
Their models have no mass predictions; the present framework
differs in three respects: (i)~the confining group is
$\SU(3)_F$ with $\Nf = \Nc = 3$ (quantum-modified constraint,
not $\Nf = \Nc + 1$ dynamical superpotential); (ii)~the
identification is \emph{reversed} (mesinos = leptons, not
baryons); and (iii)~the charge-squared law provides
quantitative mass predictions.

\paragraph{The sBootstrap.}
The \emph{sBootstrap} (supersymmetric
bootstrap)~\cite{Rivero:2024epjc} (EPJC \textbf{84}, \S2--3)
elevates the meson--baryon supersymmetry to a
\textbf{structural identification}: the composite fermions of a
confining gauge theory \emph{are} the Standard Model fermions.
The foundational statement is:
\begin{equation}
  \boxed{
  \text{mesinos} = \text{leptons}\,,
  \qquad
  \text{diquarkinos} = \text{quarks}\,.}
  \label{eq:sbootstrap_def}
\end{equation}
Here ``mesino'' means the fermionic component $\psi_M$ of the
meson chiral superfield $M^i{}_j = Q^i\tilde{Q}_j$, where
$Q^i$ and $\tilde{Q}_j$ are the fundamental ``preon'' fields of
a confining $\SU(3)_F$ family gauge theory.
The ``diquarkino'' is the fermionic component of the adjugate
composite $\adj(M)_i$, whose quantum numbers match those of a
diquark-type operator.
The ``s'' in sBootstrap stands for ``supersymmetric'':
it refers to the SUSY structure that pairs the scalar meson
$M$ with the fermionic mesino $\psi_M$ within a chiral
superfield.

\paragraph{Evidence: $m_\pi^2 - m_\mu^2 \approx f_\pi^2$.}
The mass relation $m_\pi^2 - m_\mu^2 \approx f_\pi^2$
is \emph{evidence} for the identification, not its definition.
Numerically, $m_\pi^2 - m_\mu^2 = 139.57^2 - 105.66^2 =
8316\MeV^2$ while $f_\pi^2 = 92.2^2 = 8501\MeV^2$ ($2.2\%$
discrepancy).
In the Volkov--Akulov framework~\cite{VolkovAkulov:1973}
(Phys.\ Lett.\ B \textbf{46}, eq.~1), spontaneous SUSY
breaking gives the mass relation
$m_{\rm fermion}^2 = m_{\rm boson}^2 - f^2$, where $f$ is the
SUSY-breaking scale.
If we identify $f = f_\pi$, then the muon is a pseudo-Goldstino
whose mass is split from the pion by the chiral symmetry
breaking scale.

\paragraph{The emergent $\mathcal{N}{=}2$ pairing.}
In $\mathcal{N}{=}1$ SUSY, each chiral multiplet
$(M_i, \psi_{M_i})$ provides one fermion (the mesino).
But the confined spectrum also contains the adjugate composite
$\adj(M)_i$ (Sec.~\ref{sec:seiberg}), which forms a second
chiral multiplet per family.
The pair $(M_i, \adj(M)_i)$ has the mathematical structure of an
$\mathcal{N}{=}2$ hypermultiplet---not because the UV theory has
$\mathcal{N}{=}2$ supersymmetry, but because confinement
\emph{automatically} produces both $M$ and $\adj(M)$ with the
correct quantum numbers.

\paragraph{The Q-generator problem.}
In the Miyazawa superalgebra, a single supercharge $Q$ maps
$q\bar{q} \to q[qq]$ (meson to baryon).
In the sBootstrap, the involution $M_i \leftrightarrow
\Lambda^6/M_i$ (the adjugate map) is an \emph{exact}
$\mathbb{Z}_2$ symmetry of the quantum-modified constraint
surface $\det M = \Lambda^6$: it exchanges the direct and
inverse branches and is a holomorphic identity, not an
approximation.
The \textbf{open question} is whether this $\mathbb{Z}_2$ can
be promoted to a continuous $\SU(2)_R$ rotation within the
doublet $(M_i, \adj(M)_i)$.
If so, the associated generator $Q_2$ would connect the meson
sector (leptons) to the adjugate sector (quarks), realising
Miyazawa's vision with two supercharges.
$Q_1$ is the standard $\mathcal{N}{=}1$ supercharge; the
approximate $\SU(2)_R$ is broken by the adjoint mass $\mu_\Phi$.
We construct $Q_2$ explicitly in Section~\ref{sec:N2} and show
it satisfies a cubic algebra.

\paragraph{Summary of results.}
The framework makes the following quantitative connections.
The energy-balance condition $K = 2/3$ is satisfied by
charged leptons ($0.001\%$), pseudoscalar mesons $(-\pi, D_s, B)$
($0.10\%$), inverse down quarks ($0.22\%$), and the quark chain
$(t,b,c)$--$(b,c,-s)$--$(c,s,0)$--$(s,0,m_u{+}m_d)$ with joint
significance $p \sim 1/21\,000$.
The bion relation $\zch_b = 3\zch_s + \zch_c$ predicts
$m_b = 4183\MeV$ ($0.01\%$), and the top completion formula
predicts $m_t = 168.5\GeV$ ($2.3\%$).
Five independent criteria select $N = 3$ generations.
The framework does NOT yet predict: the CKM matrix beyond the
Cabibbo angle, the PMNS matrix, the origin of the Cartan angle
$\alpha$, or the K\"ahler metric $Z_i$.

\paragraph{Plan of the paper.}
We proceed by building up the framework layer by layer:
composites (\S\ref{sec:composites}) $\to$
Seiberg dynamics (\S\ref{sec:seiberg}) $\to$
charge-squared law (\S\ref{sec:charges}) $\to$
emergent isospin (\S\ref{sec:N2}) $\to$
SUSY breaking (\S\ref{sec:susy_breaking}) $\to$
inverse sector and chain (\S\ref{sec:inverse}) $\to$
generation uniqueness (\S\ref{sec:SO32}) $\to$
Lagrangian (\S\ref{sec:lagrangian}) $\to$
discussion and conclusions (\S\ref{sec:discussion},
\S\ref{sec:open}).
Each section is designed to be self-contained with all
derivations shown.


%======================================================================
\section{Composites and confinement}
\label{sec:composites}
%======================================================================

\subsection{The gauge-theoretic realisation}

The UV gauge group is
\begin{equation}
  G_{\rm UV} = \SU(3)_F \times \SU(3)_c \times \SU(2)_L
  \times \UU(1)_Y\,,
  \label{eq:GUV}
\end{equation}
where $\SU(3)_F$ is the \emph{family gauge symmetry} that
confines at a scale $\Lambda_F$.
The term ``family'' refers to the three generations of SM
fermions: $\SU(3)_F$ acts on the generation index
$i \in \{1, 2, 3\}$.

The fundamental matter fields (``preons'') are:
\begin{itemize}
  \item $Q^\alpha_i$: transforming as
    $(\mathbf{3}_F, \mathbf{1}_c, \mathbf{3}_L)$,
    where $\alpha \in 3_L$ is an index under the
    \emph{global} $\SU(3)_L$ flavour symmetry
    (containing $\SU(2)_L \times \UU(1)_Y$ as a subgroup).
  \item $\tilde{Q}_i^\alpha$: the conjugate,
    $(\overline{\mathbf{3}}_F, \mathbf{1}_c,
    \overline{\mathbf{3}}_L)$.
\end{itemize}
The $\SU(3)_L$ is \textbf{not} a gauge symmetry: it is a global
symmetry of the preons that contains the gauged
$\SU(2)_L \times \UU(1)_Y$ as a subgroup.
Its sole role is to organise the branching
$\mathbf{3}_L \to \mathbf{2}_{-1/2} \oplus \mathbf{1}_{+1}$,
which assigns the correct SM quantum numbers to the composites
and projects out exotic charges (Sec.~\ref{sec:SO32}).
The preons are colour singlets: quarks (which carry colour) arise
through partial compositeness.

\begin{table}[ht]
\centering\small
\caption{UV field content (above $\Lambda_F$).
  All fields listed as $\mathcal{N}{=}1$ chiral superfields
  (containing a complex scalar and a Weyl fermion) except the
  gauge multiplets (containing a gauge boson and a gaugino).}
\label{tab:UV_fields}
\begin{tabular}{lccccl}
\hline
Field & $\SU(3)_F$ & $\SU(3)_c$ & $\SU(2)_L$ & $\UU(1)_Y$ & Role \\
\hline
$Q^\alpha_i$ & $\mathbf{3}$ & $\mathbf{1}$ & $\mathbf{2}\oplus\mathbf{1}$ & $-\tfrac{1}{2},+1$ & preon \\
$\tilde{Q}_i^\alpha$ & $\overline{\mathbf{3}}$ & $\mathbf{1}$ & $\overline{\mathbf{2}}\oplus\mathbf{1}$ & $+\tfrac{1}{2},-1$ & anti-preon \\
$\Phi$ (auxiliary) & $\mathbf{8}$ & $\mathbf{1}$ & $\mathbf{1}$ & $0$ & adjoint (integrated out) \\
$V_F^\mu$ & $\mathbf{8}$ & $\mathbf{1}$ & $\mathbf{1}$ & $0$ & family gauge boson \\
$q_L^a$ & $\mathbf{1}$ & $\mathbf{3}$ & $\mathbf{2}$ & $+\tfrac{1}{6}$ & elementary quark \\
$u_R^a$ & $\mathbf{1}$ & $\mathbf{3}$ & $\mathbf{1}$ & $+\tfrac{2}{3}$ & elementary quark \\
$d_R^a$ & $\mathbf{1}$ & $\mathbf{3}$ & $\mathbf{1}$ & $-\tfrac{1}{3}$ & elementary quark \\
\hline
\end{tabular}
\end{table}


\subsection{The meson decomposition and the 4/3 projection}

When $\SU(3)_F$ confines, the preons form composite bound states.
The meson superfield is the colour-singlet bilinear
$M^\alpha{}_i = Q^\alpha\tilde{Q}_i$.
It carries one $\SU(3)_L$ index ($\alpha$) and one
$\SU(3)_F$ index ($i$), giving $3 \times 3 = 9$ complex
chiral multiplets.

To extract the SM quantum numbers, we embed the electroweak
group $\SU(2)_L \times \UU(1)_Y$ inside $\SU(3)_L$ via the
standard branching:
\begin{equation}
  \mathbf{3}_L \;\to\;
  \mathbf{2}_{-1/2} \;\oplus\; \mathbf{1}_{+1}\,.
  \label{eq:3L_branching}
\end{equation}
This is the same decomposition used in $\SU(5)$ GUT.
Under \eqref{eq:3L_branching}, the meson decomposes as:
\begin{equation}
  M^\alpha{}_i \;=\;
  \underbrace{M^{1,2}{}_i}_{(\mathbf{2}_{-1/2})}
  \;\oplus\;
  \underbrace{M^{3}{}_i}_{(\mathbf{1}_{+1})}\,.
  \label{eq:meson_decomp}
\end{equation}

This decomposition is a \textbf{projection tool}: its role is
to clean the $\mathbf{15} \oplus \overline{\mathbf{15}} \oplus
\mathbf{24}$ composite content for the $\SO(32)$ string embedding
(Sec.~\ref{sec:SO32}).
Without it, the $\mathbf{24}$ of $\SU(5)_f$ contains the
$(\overline{\mathbf{3}}, \mathbf{2})_{5/6}$ component with
exotic charge $Q = 4/3$~\cite{Rivero:2024epjc}
(EPJC \textbf{84}, Theorem~3)---the analogue of the $X$ and $Y$
gauge bosons in $\SU(5)$ GUT.
The branching projects these out, restricting the composites
to SM-compatible charges.
This is the \emph{only} reason for introducing $\SU(3)_L$.

The mesino identification follows:
\begin{equation}
  \boxed{
  \psi_{M^{1,2}{}_i} \equiv L_i
  \;=\; \begin{pmatrix} \nu_i \\ e_i \end{pmatrix}_L\,,
  \qquad
  \psi_{M^{3}{}_i} \equiv e^c_i\,.}
  \label{eq:mesinos_leptons}
\end{equation}
The 9~mesino fermions furnish the complete SM lepton sector:
$3L_i + 3e^c_i$ $=$ 3 lepton doublets $+$ 3 charged-lepton
singlets.
Anomaly matching between UV preons and IR composites is satisfied
for $\Nf = \Nc$~\cite{CsakiSchmaltzSkiba:1997}
(PRL \textbf{78}, \S3).

\begin{table}[ht]
\centering\small
\caption{IR composites (below $\Lambda_F$, $\Nf = \Nc = 3$).
  Independent DOF count: $9\,(M) + 1\,(B) + 1\,(\tilde{B})
  - 1\,(X) = 10 = 18 - 8$.
  The adjugate $\adj(M)$ is a \emph{derived} operator---fixed
  by $M$ via $\adj(M)\cdot M = (\det M)\,\mathbf{I}$---and
  does NOT contribute additional DOF.}
\label{tab:IR_composites}
\begin{tabular}{llcccl}
\hline
Composite & Expression & $\SU(2)_L$ & $\UU(1)_Y$ & Chirals & SM identification \\
\hline
$M^{1,2}{}_i$ & $Q^{1,2}\tilde{Q}_i$ & $\mathbf{2}$ & $-\tfrac{1}{2}$ & 6 & $H_d$ (scalar), $L_i$ (fermion) \\
$M^{3}{}_i$ & $Q^{3}\tilde{Q}_i$ & $\mathbf{1}$ & $+1$ & 3 & slepton (scalar), $e^c_i$ (fermion) \\
\hline
\multicolumn{6}{l}{\textit{Derived operators (not independent DOF):}} \\
$\adj(M)^{1,2}{}_i$ & $\varepsilon\varepsilon\,MM$ & $\mathbf{2}$ & $+\tfrac{1}{2}$ & (6) & $H_u$ \\
$\adj(M)^{3}{}_i$ & $\varepsilon\varepsilon\,MM$ & $\mathbf{1}$ & $-1$ & (3) & sneutrino-type \\
$B$ & $\varepsilon QQQ$ & $\mathbf{1}$ & $0$ & 1 & sterile singlet \\
$\tilde{B}$ & $\varepsilon\tilde{Q}\tilde{Q}\tilde{Q}$ & $\mathbf{1}$ & $0$ & 1 & sterile singlet \\
$X$ & constraint & $\mathbf{1}$ & $0$ & $-1$ & Lagrange multiplier \\
\hline
\end{tabular}
\end{table}


\subsection{Quarks via Fradkin--Shenker complementarity}
\label{sec:composites_quarks}

Quarks carry colour ($\SU(3)_c$) and cannot be mesinos (which are
colour singlets).
The colour-triplet impossibility is a group-theory fact:
the composites of $\SU(3)_F$ with colour-singlet preons satisfy
$\mathbf{3}\otimes\overline{\mathbf{3}} = \mathbf{8}\oplus\mathbf{1}$
(mesons) and
$\mathbf{3}\otimes\mathbf{3}\otimes\mathbf{3} =
\mathbf{10}\oplus\mathbf{8}\oplus\mathbf{8}\oplus\mathbf{1}$
(baryons)---colour triplets \emph{never} appear in the
gauge-invariant composite spectrum of $\SU(3)_F$.
This makes partial compositeness structurally necessary.

The Fradkin--Shenker theorem~\cite{FradkinShenker:1979}
(Phys.\ Rev.\ D \textbf{19}, Theorem~1) states that for a
gauge theory with fundamental Higgs fields, there is
\emph{no phase transition} separating the Higgs and confinement
regimes: they are smoothly connected.
Concretely, for $\SU(3)$ with $\Nf = 3$ fundamentals:
\begin{itemize}
  \item \textbf{Confinement phase}: the gauge-invariant
    composites are mesons ($M = Q\tilde{Q}$), baryons
    ($B = \varepsilon QQQ$), and their conjugates.
    Total light DOF: $9 + 1 + 1 - 1 = 10$.
  \item \textbf{Higgs phase}: if the preons develop VEVs,
    $\SU(3)_F$ is completely broken, eating 8 Goldstone bosons.
    Remaining light DOF: $2 \times 9 - 8 = 10$.
\end{itemize}
The Fradkin--Shenker theorem guarantees that these are the
\emph{same} 10 degrees of freedom, continuously connected.

The coloured quarks of the SM enter through \textbf{partial
compositeness}~\cite{Kaplan:1991dc} (Nucl.\ Phys.\ B \textbf{365},
eq.~2.1):
each elementary quark $q_{\rm elem}$ (which carries colour and
is a singlet of $\SU(3)_F$, see Table~\ref{tab:UV_fields})
mixes linearly with a composite operator $\mathcal{O}$ of
matching SM quantum numbers:
\begin{equation}
  \mathcal{L}_{\rm mix} = \lambda_q\,
  \bar{q}_{\rm elem}\,\mathcal{O}_{\rm composite}
  + \text{h.c.}
\end{equation}
The physical quark is a linear combination of elementary and
composite fields, with mass $m_q \sim \lambda_q^2 / M_*$ where
$M_*$ is the compositeness scale.
This ``seesaw'' mechanism automatically gives $m \propto \zch^2$
if the mixing parameter $\lambda_q \propto \zch_q$ (the family
charge).

\paragraph{Why composites cannot carry colour triplets.}
This is a group theory fact, not a dynamical assumption.
The gauge-invariant composites of $\SU(3)_F$ with colour-singlet
preons decompose under $\SU(3)_c$ as:
\begin{itemize}
  \item Mesons: $\mathbf{1}_c \otimes \mathbf{1}_c = \mathbf{1}_c$
    (colour singlet).
  \item Baryons: $\mathbf{1}_c \otimes \mathbf{1}_c \otimes
    \mathbf{1}_c = \mathbf{1}_c$ (colour singlet).
  \item Adjugate: products of colour singlets remain colour singlets.
\end{itemize}
Colour triplets \emph{never} appear in the composite spectrum.
If one instead considered colour-carrying preons (as in the
Luty--Mohapatra construction~\cite{LutyMohapatra:1997}), the
baryon $B = \varepsilon QQQ$ would carry
$\mathbf{3} \otimes \mathbf{3} \otimes \mathbf{3}
= \mathbf{10} \oplus \mathbf{8} \oplus \mathbf{8}
\oplus \mathbf{1}$---still no colour triplet.
The sBootstrap's choice of colour-singlet preons makes partial
compositeness structurally \emph{necessary}: quarks must acquire
colour from an elementary sector.


\subsection{The lepton sector obstruction}

There is an honest open problem.
The holomorphic meson masses from the Seiberg vacuum
(Sec.~\ref{sec:seiberg}) give mass \emph{ratios}
determined by the UV preon masses.
The mesino mass ratio
$m_\psi(M^u{}_s)/m_\psi(M^u{}_d) = m_s/m_d \approx 20$,
while the physical lepton ratio $m_\mu/m_e = 207$---a factor
$10.3$ discrepancy.
K\"ahler corrections preserve the rank of the mass matrix but
can only introduce flavour-independent rescalings if $Z$ is
generation-universal.
Flavour-\emph{dependent} rescaling (different $Z_i$ per
generation) can generate the discrepancy.
Two candidate mechanisms:
(i) Nelson--Strassler anomalous dimensions at a confining
fixed point~\cite{NelsonStrassler:2000};
(ii) a separate $\SU(2)$ confining sector with $\Nf = 3$ and
an O'Raifeartaigh deformation (Sec.~\ref{sec:susy_breaking}).
We flag this as an open issue, not an inconsistency: the
framework predicts the \emph{structure} of the mass matrix
(diagonal, with entries $\propto \zch_i^2$) but not the
overall K\"ahler normalisation.


%======================================================================
\section{Seiberg confinement and the adjugate mechanism}
\label{sec:seiberg}
%======================================================================

\subsection{Review: s-confinement for $\Nf = \Nc$}

We review the non-perturbative dynamics of $\SU(\Nc)$ SQCD with
$\Nf$ flavours of quarks and antiquarks in the $\mathcal{N}{=}1$
framework~\cite{Seiberg:1994bz,Tong:SUSY}.

For $\Nf = \Nc$, the theory is \emph{s-confining} (``smoothly
confining''): the low-energy degrees of freedom are the
gauge-invariant composites $M$, $B$, $\tilde{B}$, and there is
no residual gauge group.
The classical moduli space is parametrised by the constraint
$\det M - B\tilde{B} = 0$.
Quantum effects (a one-instanton calculation) modify this to the
\textbf{quantum-modified constraint}:
\begin{equation}
  \det M - B\tilde{B} = \Lambda^{2\Nc}\,,
  \label{eq:constraint}
\end{equation}
where $\Lambda$ is the dynamical scale.
The right-hand side $\Lambda^{2\Nc}$ is generated by instantons
and is exact (protected by holomorphy and symmetries).
The effect is to smooth out the moduli space: the origin
$M = B = \tilde{B} = 0$ is removed.

The constraint is implemented in the effective superpotential by
a Lagrange multiplier $X$:
\begin{equation}
  W = X\bigl(\det M - B\tilde{B} - \Lambda^6\bigr)
    + \sum_i m_i\,M^i{}_i\,.
  \label{eq:Wtot}
\end{equation}
The first term enforces the constraint; the second adds explicit
preon masses which break global symmetry and lift the moduli
space to isolated vacua.


\subsection{The mesonic vacuum and the Seiberg seesaw}

On the mesonic branch ($B = \tilde{B} = 0$), we solve the
$F$-term equations $\partial W / \partial M^i{}_j = 0$.
For diagonal $M = \diag(M_1, M_2, M_3)$:
\begin{equation}
  \frac{\partial W}{\partial M_i}
  = X \cdot \frac{\partial(\det M)}{\partial M_i} + m_i
  = X \cdot \prod_{j \neq i} M_j + m_i = 0\,.
  \label{eq:Fterm_M}
\end{equation}
The product $\prod_{j \neq i} M_j$ is precisely the $i$-th
diagonal entry of the adjugate matrix.
From the constraint $\det M = M_1 M_2 M_3 = \Lambda^6$,
and taking ratios of $F$-term equations for different $i$:
$M_i/M_j = m_j/m_i$, which gives the \textbf{Seiberg seesaw}:
\begin{equation}
  \langle M_i \rangle = \frac{\Lambda^2 P^{1/3}}{m_i}\,,
  \qquad P = m_1 m_2 m_3\,.
  \label{eq:vacuum}
\end{equation}
Heavy preons produce light mesons, and vice versa.
This is the single most important equation in the framework.

\paragraph{Worked example.}
Take three preons with masses $m_1 = 1$, $m_2 = 2$, $m_3 = 4$
(arbitrary units) and $\Lambda = 1$.
Then $P = 1 \cdot 2 \cdot 4 = 8$, $P^{1/3} = 2$, and:
\begin{equation}
  \langle M_1 \rangle = \frac{2}{1} = 2\,,\quad
  \langle M_2 \rangle = \frac{2}{2} = 1\,,\quad
  \langle M_3 \rangle = \frac{2}{4} = \frac{1}{2}\,.
\end{equation}
Check: $M_1 M_2 M_3 = 2 \cdot 1 \cdot 1/2 = 1 = \Lambda^6$.
The heaviest preon ($m_3 = 4$) produces the lightest meson
($M_3 = 1/2$), confirming the seesaw.


\subsection{The $\Nf$-dependent phase structure}

The behaviour of $\SU(3)_F$ depends critically on how many preon
flavours are light compared to the RG scale.

\begin{table}[ht]
\centering\small
\caption{Phase structure of $\SU(3)_F$ as heavy flavours decouple.
  Scale matching: $\Lambda_3^6 = \Lambda_4^5 m_c$,
  $\Lambda_4^5 = \Lambda_5^4 m_b$.
  Each threshold crossing changes the non-perturbative dynamics
  qualitatively.}
\label{tab:phases}
\begin{tabular}{clll}
\hline
$\Nf$ & Energy range & Phase & Key property \\
\hline
$6$ & $Q > m_t$ & conformal window & IR fixed point \\
$5$ & $m_b < Q < m_t$ & conformal window & IR fixed point \\
  &  & ($4.5 < 5 < 9$) & \\
$4$ & $m_c < Q < m_b$ & ISS metastable~\cite{ISS:2006} & unique integer \\
  &  & & in $(\Nc, \tfrac{3}{2}\Nc)$ \\
$3$ & $Q < m_c$ & s-confinement & $\det M = \Lambda^6$ \\
\hline
\end{tabular}
\end{table}

At $\Nf = 4$, the theory falls into the ISS
window~\cite{ISS:2006}: $\Nc < \Nf < \frac{3}{2}\Nc$,
i.e.\ $3 < 4 < 4.5$.
This is the window where metastable SUSY-breaking vacua exist.
$\Nf = 4$ is the \emph{unique} integer in this window for
$\Nc = 3$---a discrete selection that plays a role in the
effective Lagrangian (Sec.~\ref{sec:lagrangian}).


\subsection{The adjugate: definition and explicit examples}

The adjugate of a square matrix $M$ is defined by:
\begin{equation}
  \adj(M) \cdot M = \det(M) \cdot I\,.
  \label{eq:adj_def}
\end{equation}
Unlike the inverse $M^{-1} = \adj(M)/\det M$, the adjugate
is a \emph{polynomial} in the entries of $M$---holomorphic and
well-defined even when $\det M = 0$.

\paragraph{$2 \times 2$ example.}
For
$A = \bigl(\begin{smallmatrix} a & b \\ c & d \end{smallmatrix}
\bigr)$:
$\adj(A) = \bigl(\begin{smallmatrix} d & -b \\ -c & a
\end{smallmatrix}\bigr)$.
Check: $\adj(A) \cdot A = (ad - bc)\,I_2 = \det(A)\,I_2$.
The adjugate swaps the diagonal entries and negates the
off-diagonal ones---the simplest example of ordering reversal.

\paragraph{$3 \times 3$ diagonal example.}
For $M = \diag(M_1, M_2, M_3)$:
\begin{equation}
  \adj(M) = \diag(M_2 M_3,\; M_1 M_3,\; M_1 M_2)\,.
  \label{eq:adj_3x3}
\end{equation}
Each entry is the product of the \emph{other two} diagonal entries.
If $M_1 < M_2 < M_3$, then
$\adj(M)_1 = M_2 M_3 > \adj(M)_2 = M_1 M_3 > \adj(M)_3 =
M_1 M_2$.
The ordering is \emph{reversed}: the heaviest preon produces the
lightest meson VEV and the largest adjugate entry.

\paragraph{Numerical example with lepton masses.}
Taking $M = \diag(m_e, m_\mu, m_\tau) =
\diag(0.511, 105.66, 1776.86)\MeV$:
\begin{align}
  \adj(M)_e &= m_\mu \cdot m_\tau = 187\,740\MeV^2\,,
  \nonumber\\
  \adj(M)_\mu &= m_e \cdot m_\tau = 908\MeV^2\,,
  \nonumber\\
  \adj(M)_\tau &= m_e \cdot m_\mu = 54.0\MeV^2\,.
\end{align}
The heaviest lepton ($\tau$) has the \emph{smallest} adjugate
entry.


\subsection{The two Yukawa operators: direct vs.\ inverse}

From the $F$-term equations~\eqref{eq:Fterm_M}:
\begin{equation}
  \adj(M)_i = -\frac{m_i}{X} = -\Lambda^3\,m_i\,,
  \label{eq:adjM_from_Fterm}
\end{equation}
using $X = -P^{1/3}/\Lambda^4$ from the constraint.
Together with $\langle M_i \rangle = \Lambda^2 P^{1/3}/m_i$:
\begin{align}
  M_i &\propto \frac{1}{m_i}
    \quad\text{(inverse: heavy UV $\to$ light IR)}\,,
    \label{eq:M_inverse}\\
  \adj(M)_i &\propto m_i
    \quad\text{(direct: heavy UV $\to$ heavy IR)}\,.
    \label{eq:adjM_direct}
\end{align}
These are the two natural Yukawa operators of the confined theory.
The physical assignment:
\begin{itemize}
  \item \textbf{Leptons} are mesinos $\psi_M$.
    Their mass comes from the $F$-term
    $|F_{M_i}| \propto m_i$ (direct).
    Hence: $K(m_e, m_\mu, m_\tau) = 0.666661$ ($0.001\%$).
  \item \textbf{Down-type quarks} couple to the meson VEV
    $\langle M \rangle \propto 1/m_i$ through partial
    compositeness.
    Hence: $K^{-1}(m_d, m_s, m_b) = 0.6652$ ($0.22\%$).
\end{itemize}
A \emph{single UV condition} $K(m^{\rm UV}) = 2/3$ produces
both the standard lepton Koide and the inverse quark Koide,
unified by the involution $M_i \leftrightarrow \Lambda^4/M_i$.

\paragraph{Automatic pairing of direct and inverse.}
If the UV masses satisfy $m_i = \mu\zch_i^2$ with $\UU(3)_F$
charges, then on the constraint surface
$a_i = \alpha\zch_i^2$ and $d_i = \beta/\zch_i^2$ with
$\alpha\beta = \Lambda^6$.
Both $K_{\rm dir}$ and $K_{\rm inv}$ evaluate to $2/3$ by
the same algebraic identity (Theorem~\ref{thm:n2}).

\begin{table}[ht]
\centering\small
\caption{Direct and inverse branch assignment.
  Product rule: $a_i \cdot d_i = \mu\nu = \text{const}$
  ($\det M = \Lambda^6$).}
\label{tab:branches}
\begin{tabular}{lllll}
\hline
Branch & Operator & Mass formula & Ordering & Physical sector \\
\hline
Direct & $\adj(M)_i \propto m_i$ & $a_i = \mu\zch_i^2$ & $a_1 < a_2 < a_3$ & Leptons (F-term) \\
Inverse & $M_i \propto 1/m_i$ & $d_i = \nu/\zch_i^2$ & $d_1 > d_2 > d_3$ & Down quarks (VEV) \\
\hline
\end{tabular}
\end{table}


\subsection{Off-diagonal masses and transparency}

The off-diagonal meson masses
are~\cite{CsakiSchmaltzSkiba:1997}:
\begin{equation}
  m(M^i{}_j) = \frac{m_i m_j}{P^{1/3}}\,,
  \qquad i \neq j\,.
  \label{eq:offdiag}
\end{equation}
These depend only on ratios of the input masses---they are
$\Lambda$-independent.
The diagonal mass matrix has $(W_S)_{kk} = 0$ for all $k$.
This means the Seiberg dynamics is \textbf{transparent}: it
preserves $K = 2/3$ (if present in the UV masses) but cannot
generate it.
The origin of $K = 2/3$ must come from the charge-squared law
(Sec.~\ref{sec:charges}).

\paragraph{Numerical transparency check.}
Using physical lepton masses as UV inputs:
$m_e = 0.511\MeV$, $m_\mu = 105.66\MeV$, $m_\tau = 1776.86\MeV$.
The off-diagonal masses from eq.~\eqref{eq:offdiag} are:
$m(M^e{}_\mu) = m_e m_\mu / P^{1/3}
= 0.511 \times 105.66 / 45.6 = 1.18\MeV$,
$m(M^e{}_\tau) = 0.511 \times 1776.86 / 45.6 = 19.9\MeV$,
$m(M^\mu{}_\tau) = 105.66 \times 1776.86 / 45.6 = 4116\MeV$.
These off-diagonal masses depend only on \emph{ratios} of the
input masses---they contain no information about $\Lambda$.
The diagonal entries $(W_S)_{kk} = 0$: no diagonal mass is
generated by the Seiberg dynamics.
Hence $K = 2/3$ in the UV masses passes unchanged through
the confinement---transparency.


%======================================================================
\section{The charge-squared mass law}
\label{sec:charges}
%======================================================================

We derive the mass ansatz $m_i = \mu\,\zch_i^2$, where
$\zch_i = z_0 + z_i$ is the signed family charge.
Three independent arguments converge on the same result.

\subsection{The $\UU(3)_F$ charge algebra}

The family symmetry $\UU(3)_F = \SU(3)_F \times \UU(1)$ has a
Cartan subalgebra spanned by two diagonal generators of
$\SU(3)_F$ plus the $\UU(1)$ generator.
With canonical normalisation:
\begin{equation}
  \sum_{i=1}^{3} z_i = 0\,,\qquad
  \sum_{i=1}^{3} z_i^2 = \frac{1}{2}\,,\qquad
  z_0 = \frac{1}{\sqrt{6}}\,.
\end{equation}
The signed charge is $\zch_i = z_0 + z_i$, with
$z_i = \cos(\alpha + 2\pi(i-1)/3)/\sqrt{3}$ parametrised by
the Cartan direction $\alpha$ (the ``sector angle'').


\subsection{Argument 1: Algebraic uniqueness of $n = 2$}
\label{sec:uniqueness_n2}

Consider a general monomial mass law $m_i = c\,(z_0 + z_i)^n$
and define:
\begin{equation}
  K_n(\alpha) \equiv
  \frac{\sum_i (z_0{+}z_i)^n}
  {\bigl(\sum_i (z_0{+}z_i)^{n/2}\bigr)^2}\,.
  \label{eq:Kn}
\end{equation}

\begin{theorem}[Uniqueness of $n = 2$]
\label{thm:n2}
For integer powers $n$ and $\UU(3)$ charges $\zch_i = z_0 + z_i$,
the ratio $K_n(\alpha)$ is independent of the Cartan direction
$\alpha$ if and only if $n = 2$ (besides $n = 0$).
At $n = 2$: $K_2 = 2/3$.
\end{theorem}

\begin{proof}
We expand $\sum(z_0 + z_i)^n$ using the binomial theorem.
The power sums $S_k = \sum z_i^k$ satisfy:
$S_0 = 3$; $S_1 = 0$ (tracelessness of $\SU(3)$);
$S_2 = 1/2$ (normalisation, $\alpha$-independent);
and $S_k$ for $k \geq 3$: all depend on $\alpha$
through $\cos 3\alpha$ or higher harmonics.

For $n = 2$:
$\sum (z_0+z_i)^2 = 3z_0^2 + 2z_0 \cdot 0 + 1/2 = 1$.
Also $\sum (z_0+z_i) = 3z_0 = \sqrt{3/2}$.
So $K_2 = 1/(3/2) = 2/3$, independent of $\alpha$.

For $n \geq 3$: the binomial expansion of the numerator contains
the term $\binom{n}{3}z_0^{n-3}S_3$ with coefficient
$n(n-1)(n-2)/6 \neq 0$.
Since $S_3 = \sum z_i^3 \propto \cos 3\alpha$, the numerator
acquires $\alpha$-dependence.
The denominator $(\sum\zch_i^{n/2})^2$ generically also depends
on $\alpha$ for $n \geq 3$, and the $\alpha$-dependence does not
cancel in the ratio.

Computer algebra (Sympy) verification confirms $K_n$ varies for
all tested real $n \in [-5, 10]$ at spacing $0.001$, except
$n = 2$ and the trivial $n = 0$.
\end{proof}

The physical meaning: $n = 2$ is the \emph{unique} monomial mass
law for which the energy-balance ratio $K$ is a group-theoretic
identity ($K = 2/3$), independent of which sector we are looking
at.
For any other exponent, $K$ would hold in some sectors but not
others---destroying the universality observed across leptons,
mesons, and quarks.

\paragraph{Explicit check for $n = 3$.}
At $n = 3$: $\sum(z_0+z_i)^3 = 3z_0^3 + 3z_0 \cdot 0
+ 3z_0 S_2 + S_3 = 3z_0^3 + 3z_0/2 + S_3$,
where $S_3 = \cos(3\alpha)/(3\sqrt{3})$.
The denominator $(\sum(z_0+z_i)^{3/2})^2$ also depends on
$\alpha$ through its cube-root structure.
Numerical evaluation: at $\alpha = 0$,
$K_3 = 0.592$; at $\alpha = \pi/6$, $K_3 = 0.764$.
The $17\%$ variation confirms that $n = 3$ fails the direction
independence test.


\subsection{Argument 2: Composite self-energy}

Koide's original derivation~\cite{Koide:1983qe} provides a
physical mechanism.
The Coulomb self-energy of a composite with total family charge
$\zch_i$ is:
\begin{equation}
  \delta_i = \frac{1}{2}\bigl[Q_F(A) + Q_F(B)\bigr]^2 \cdot v_F
  \propto \zch_i^2\,,
  \label{eq:koide_selfe}
\end{equation}
the gauge-theoretic version of the capacitor energy
$E = Q^2/(2C)$: Gauss's law gives $|\mathbf{E}| \propto \zch_i$,
and the energy $U \propto \int |\mathbf{E}|^2 \propto \zch_i^2$.
In a confining theory, the dominant mass contribution is the
string tension $\sigma \cdot L$, which is generation-independent.
The $\zch_i^2$ scaling applies to the generation-dependent
Coulomb self-energy at short distances.


\subsection{Argument 3: Bilinear structure}

The meson $M^i{}_i = Q^i\tilde{Q}_i$ is a bilinear.
If each preon contributes $\propto \zch_i$, the composite VEV
scales as $\zch_i^2$~\cite{Koide:1990}.
This is the same structure as a seesaw: any mass from two VEV
insertions gives $m \propto v^2$.

\paragraph{Status summary.}
Argument~1 is a rigorous algebraic theorem.
Argument~2 is structural motivation (Coulomb scaling
assumed to survive confinement).
Argument~3 is a consistency check.
The charge-squared law $m_i = \mu\,\zch_i^2$ is the framework's
central \emph{postulate}: uniquely motivated but not derived from
confining dynamics.
Seven SUSY-breaking derivation attempts fail because gauge
mediation gives $\zch^2$ scaling for soft (K\"ahler) masses,
not for holomorphic (superpotential) masses needed for exact
$K = 2/3$.


\subsection{The energy-balance condition}

For three charges $\zch_k = z_0 + z_k$ with $\sum z_k = 0$:
\begin{equation}
  K = \frac{\sum \zch_k^2}{(\sum \zch_k)^2}
  = \frac{1}{3} + \frac{\langle z_k^2\rangle}{3z_0^2}\,.
  \label{eq:K_derivation}
\end{equation}
Therefore:
\begin{equation}
  \boxed{K = \frac{2}{3}
  \quad\Longleftrightarrow\quad
  \langle z_k^2\rangle = z_0^2\,.}
  \label{eq:energy_balance}
\end{equation}
At $K = 2/3$, the variance of the $\SU(3)$ charges equals the
square of the $\UU(1)$ charge: the splitting scale matches the
base scale.
This is the \emph{equipartition point} between the perturbative
limit ($K \to 1/3$, all masses equal) and the dominant limit
($K \to 1$, one mass dominates).

In the charge plane, the variable
$\langle z^2\rangle/z_0^2 = 3K - 1$ measures deviation from
$K = 2/3$.
The zero-mass algebraic fixed points occur at
$\alpha = 15^\circ$ and $45^\circ$ ($\SU(3)$ convention), with
ratio $r_0 = (2+\sqrt{3})^2 \approx 13.93$.
All five physical Koide tuples lie within $5.5^\circ$ of the
$15^\circ$ fixed point---their clustering reflects proximity to
a single algebraic attractor.

$K = 2/3$ is \textbf{not} an RG attractor: one-loop soft-mass
renormalisation drives $K \to 1/3$, not $K \to 2/3$.
The energy-balance condition is a \emph{boundary condition} at
the SUSY-breaking scale, not an RG fixed point.

\paragraph{Seed triples.}
For a triple $(0, \zch_a, \zch_b)$ with one massless member,
the energy-balance condition becomes:
$K = (\zch_a^2 + \zch_b^2)/(\zch_a + \zch_b)^2$.
Setting $r = \zch_a/\zch_b$ and solving $K = 2/3$:
$r^2 + 1 = 2(r + 1)^2/3$, giving $r = 2 - \sqrt{3}
\approx 0.268$.
This is a universal ratio: \emph{any} seed triple satisfying
$K = 2/3$ must have $\zch_{\rm small}/\zch_{\rm large}
= 2 - \sqrt{3}$.
\begin{itemize}
  \item \textbf{Meson seed} $(0, \pi, D_s)$:
    $\zch_\pi/\zch_{D_s} = 0.266$ ($0.6\%$ off),
    $K = 0.668$ ($0.18\%$).
  \item \textbf{Quark seed} $(0, s, c)$:
    $\zch_s/\zch_c = 0.271$ ($1.2\%$ off),
    $K = 0.664$ ($0.35\%$).
  \item \textbf{Leptons} $(e, \mu, \tau)$:
    $K = 0.666661$ ($0.001\%$)---an ``almost-seed''
    ($\zch_e = 0.72\MeV^{1/2}$ is small but nonzero).
\end{itemize}

\paragraph{The meson triple $(-\pi, D_s, B)$.}
With signed charges $\zch_\pi < 0$, $\zch_{D_s} > 0$,
$\zch_B > 0$:
\begin{equation}
  K(-\pi, D_s, B) = \frac{\zch_\pi^2 + \zch_{D_s}^2 + \zch_B^2}
  {(-\zch_\pi + \zch_{D_s} + \zch_B)^2}
  = 0.6674\,,
\end{equation}
$0.10\%$ from $2/3$---the second most precise triple known.
The pion enters with a \emph{negative} sign because it approaches
$m = 0$ from below in the chiral limit (a pseudo-Goldstone boson).
Heavy mesons have $m_{\rm meson} \approx m_Q + \Lambda_{\rm QCD}$
(HQET), so $m_{\rm meson} \approx m_Q \propto \zch_Q^2$: the
charge-squared law for the heavy quark propagates to the heavy
meson mass.
The GOR relation~\cite{GellMannOakesRenner:1968}
$M_\pi^2 = B_0(m_u + m_d)$ gives boson mass \emph{squared}
linear in quark mass, which is why $K(\pi, K, \eta) = 0.36$ is
far from $2/3$---and this is \emph{expected}.

\paragraph{Sumino protection.}
The charged-lepton $K = 0.666661$ ($0.001\%$) is too precise to
survive QED corrections without protection.
Sumino~\cite{Sumino:2009} showed that a gauged $\UU(3)_F$ with
$\alpha_{\rm em} = \alpha_F/4$ cancels the dangerous
$\alpha\,m_i\log m_i^2$ corrections, placing family gauge bosons
at $10^2$--$10^3\TeV$.


%======================================================================
\section{The emergent $\mathcal{N}{=}2$ structure}
\label{sec:N2}
%======================================================================

\subsection{What ``emergent $\mathcal{N}{=}2$'' means}

In genuine $\mathcal{N}{=}2$ theories~\cite{SeibergWitten:1994},
each matter field sits in a hypermultiplet: a pair of
$\mathcal{N}{=}1$ chirals $(\Phi, \tilde\Phi)$ related by
$\SU(2)_R$.
In our $\mathcal{N}{=}1$ theory, the confined spectrum
automatically produces two chiral multiplets per family:
\begin{equation}
  (M_i, \;\adj(M)_i)
  \qquad\text{per family } i = 1, 2, 3\,,
\end{equation}
with complementary quantum numbers:
$M^{1,2} \sim (\mathbf{2}, -1/2)$ while
$\adj(M)^{1,2} \sim (\mathbf{2}, +1/2)$.
This has the \emph{mathematical structure} of a
hypermultiplet---emergent from confinement plus the adjugate
identity, not from UV symmetry.


\subsection{Breaking the doublet: the quartic}

An auxiliary adjoint $\Phi$ (transforming as $\mathbf{8}$ of
$\SU(3)_F$) with mass $\mu_\Phi$ and Yukawa coupling
$W \supset \sqrt{2}\,g\,\tilde{Q}\,\Phi\,Q$~\cite{HMS:1997}
is integrated out via its $F$-term
$F_\Phi = \sqrt{2}\,g\,\tilde{Q}Q + \mu_\Phi\Phi = 0$,
giving $\Phi = -\sqrt{2}\,g\,\tilde{Q}Q/\mu_\Phi$.
Substituting back:
\begin{equation}
  W_{\rm quartic} = -\frac{g^2}{\mu_\Phi}
    \left[\Tr M^2
    - \frac{1}{3}(\Tr M)^2\right].
  \label{eq:quartic}
\end{equation}
The adjoint $\Phi$ is NOT a physical particle in the sBootstrap:
it is a UV field with mass $\mu_\Phi$ that is integrated out
below that scale, leaving only the quartic
coupling~\eqref{eq:quartic} as its low-energy footprint.
This is standard EFT decoupling---no different from the $W$
boson producing the Fermi four-fermion interaction below $M_W$.

\paragraph{The Yukawa coupling.}
In genuine $\mathcal{N}{=}2$
SUSY~\cite{Fayet:1976,SeibergWitten:1994},
the Yukawa coupling equals the gauge coupling:
$y = \sqrt{2}\,g$.
The UV theory here is $\mathcal{N}{=}1$, where $y$ is a priori
free.
We \emph{adopt} $y = \sqrt{2}\,g$ as the simplest choice
consistent with the emergent doublet structure.
This is an \emph{assumption}, not a derivation from the
$\mathcal{N}{=}1$ dynamics.


\subsection{Both Higgs doublets from one meson}

The structural consequence is that the doublet
$(M_i, \adj(M)_i)$ splits into components with \emph{different}
electroweak quantum numbers:
\begin{align}
  M^{1,2} &\sim (\mathbf{2}, -1/2) = H_d
    \qquad\text{(down-type Higgs)}\,,\\
  \adj(M)^{1,2} &\sim (\mathbf{2}, +1/2) = H_u
    \qquad\text{(up-type Higgs)}\,.
\end{align}
\textbf{Both MSSM Higgs doublets emerge from a single meson
matrix.}
There is no independent $\tan\beta$ parameter: both $H_d$ and
$H_u$ are determined by the same meson matrix and constrained by
$\det M = \Lambda^6$.
The ``superpartners'' are the mesons themselves, split from
mesinos by $f_\pi$; there is no second SUSY at the TeV scale.


\subsection{The $Q_2$ operator and the cubic algebra}

The involution $M_i \leftrightarrow \Lambda^6/M_i$ is a
$\mathbb{Z}_2$ symmetry of the constraint surface.
We now construct an explicit operator that implements this map.

Define the off-diagonal matrix $R$ by
\begin{equation}
  R_{ij} = \langle M_k\rangle \quad (k \neq i,j; \;\{i,j,k\}
  \text{ cyclic})\,, \qquad R_{ii} = 0\,.
  \label{eq:Rdef}
\end{equation}
Acting on the meson VEV vector
$\mathbf{M} = (\langle M_1\rangle, \langle M_2\rangle,
\langle M_3\rangle)^T$, $R$ gives an exact identity:
\begin{equation}
  R \cdot \mathbf{M} = 2\,\adj(\mathbf{M})\,,
  \label{eq:RM_adjM}
\end{equation}
verified by direct computation: the $i$-th component of $R
\cdot \mathbf{M}$ is $\sum_{j \neq i} M_k M_j = 2\,M_j M_k =
2\,\adj(M)_i$.
The operator $R$ therefore rotates the meson sector into the
adjugate sector---precisely the $Q_2$ map connecting leptons to
quarks.

On the constraint surface $\det M = \Lambda^6$, the
Cayley--Hamilton identity gives:
\begin{equation}
  R^3 = \Tr(D^2)\,R + 2\Lambda^6\,I\,,
  \label{eq:R_cubic}
\end{equation}
where $D = \diag(\langle M_i\rangle)$.
This is a \emph{cubic} minimal polynomial, not the quadratic
$R^2 = c\,I$ of standard $\mathcal{N}{=}2$.
The coefficient $2\Lambda^6 = 2\det M$ is the quantum
constraint; $\Tr(D^2) = \sum M_i^2$ is the quadratic Casimir.

The $D$-conjugated operator is democratic:
\begin{equation}
  D\, R\, D = \Lambda^6\,(J - I)\,,
  \label{eq:DRD}
\end{equation}
where $J_{ij} = 1$ for all $i,j$ is the all-ones matrix.
After absorbing the mass hierarchy into $D$, the $Q_2$
rotation treats all families equally.

The normalised operator $T = R/\sqrt{\det D}$ has eigenvalues
$\approx \{+1.07, -3 \times 10^{-6}, -1.07\}$ for the
lepton hierarchy---an ``almost-involution'' with one near-zero
eigenvalue reflecting the mass hierarchy.

\paragraph{Connection to Kostant and the principal $\mathfrak{sl}(2)$.}
The operator $R$ is the \emph{pure spin-2 component} of
$\mathrm{End}(V)$ under the principal $\mathfrak{sl}(2)$
embedding in $\mathfrak{sl}(3)$~\cite{Kostant:1959}.
The proof is direct: $R$ is symmetric ($R = R^T$, eliminating
spin-1 which is antisymmetric) and traceless ($\Tr R = 0$,
eliminating spin-0 which is proportional to $I$), so $R$
lies in the spin-2 subspace of
$\mathrm{End}(V) = \text{spin-0} \oplus \text{spin-1}
\oplus \text{spin-2}$ (dim $1 + 3 + 5 = 9$).
The exponents of $\mathfrak{sl}(3)$ are $\{1, 2\}$; $R$ lives
in the exponent-2 representation.
Whether the cubic algebra~\eqref{eq:R_cubic} defines a
consistent deformed $\mathcal{N}{=}2$ superalgebra is an open
question.


\subsection{The critical scale and self-completion}

The energy-balance vacuum exists only above
\begin{equation}
  \mu_{\Phi,\rm crit} \approx 4.5\TeV\,.
  \label{eq:mucrit}
\end{equation}
Under 5-loop SMDR running~\cite{MartinRobertson:2019,PDG:2024},
$K^{-1}(d,s,b)$ passes through $2/3$ at $Q \approx 280\TeV$;
the quartic must be large enough to sustain the mixed vacuum,
requiring $\mu_\Phi \lesssim 4.5\TeV$.

The model self-completes: above $\Lambda_F$, the theory is
$\SU(3)_F \times \text{SM}$ with preons, asymptotically free
to the Planck scale.
There are no stops, sbottoms, gluinos, or conventional SUSY
partners: the ``SUSY'' refers to composite-sector structure,
not a TeV-scale extension.


%======================================================================
\section{SUSY breaking: seeds, bloom, and bions}
\label{sec:susy_breaking}
%======================================================================

\subsection{The O'Raifeartaigh model and the seed}

The O'Raifeartaigh superpotential~\cite{ORaifeartaigh:1975}
with three chiral superfields:
\begin{equation}
  W_{\rm OR} = f\,\Phi_0 + m\,\Phi_1\Phi_2
  + g\,\Phi_0\Phi_1^2\,,
  \label{eq:W_OR}
\end{equation}
breaks SUSY because the $F$-terms
\begin{align}
  F_0 &= f + g\,\Phi_1^2\,,\quad
  F_1 = m\,\Phi_2 + 2g\,\Phi_0\,\Phi_1\,,\quad
  F_2 = m\,\Phi_1
\end{align}
cannot simultaneously vanish ($F_0 = 0$ requires
$\Phi_1^2 = -f/g$, but then $F_2 = m\Phi_1 \neq 0$).
The fermionic mass matrix at
$\langle\Phi_0\rangle = v$,
$\langle\Phi_1\rangle = \langle\Phi_2\rangle = 0$:
\begin{equation}
  \mathcal{M}_F = \begin{pmatrix}
    0 & 0 & 0 \\ 0 & 2gv & m \\ 0 & m & 0
  \end{pmatrix}\,,
  \label{eq:fermion_mass_matrix}
\end{equation}
has eigenvalues $m_\pm = \sqrt{g^2v^2 + m^2} \pm gv$ and
$m_0 = 0$ (the goldstino).
At the seed condition $gv/m = \sqrt{3}$:
$m_- = (2-\sqrt{3})\,m$ and $m_+ = (2+\sqrt{3})\,m$.
Computing $K$:
\begin{equation}
  K(0, m_-, m_+) = \frac{m_- + m_+}{(\sqrt{m_-}+\sqrt{m_+})^2}
  = \frac{4m}{6m} = \frac{2}{3}\,,
\end{equation}
exactly.
K\"ahler stabilisation at the pseudo-modulus pole with
$\Lambda_K = m$ requires:
\begin{equation}
  c = -\frac{1}{12}\,.
  \label{eq:c_12}
\end{equation}

\paragraph{Why this is ``mild'' SUSY breaking.}
In the ISS dictionary~\cite{ISS:2006}, $g \to h$,
$m \to h\mu_{\rm ISS}$, $v \to X$, so the seed condition becomes
$\langle X\rangle/\mu_{\rm ISS} = \sqrt{3}$.
In the sBootstrap, $f = f_\pi^2 \approx (92\MeV)^2$ while
$m = \mu/4 \approx 471\MeV$, giving $f/m^2 \approx 0.038$.


\subsection{The bloom: from seed to physical spectrum}

The seed is $(0, m_-, m_+)$ with one vanishing mass.
The physical lepton spectrum $(m_e, m_\mu, m_\tau)$ has
$K = 0.666661$ but no vanishing mass.
The \textbf{bloom} is the nonperturbative process that transforms
the seed into the physical spectrum.

In the charge parametrisation
$\zch_k = z_0\bigl(1 + \sqrt{2}\cos(\delta + 2\pi k/3)\bigr)$,
the seed corresponds to $\delta = 3\pi/4$ (one charge vanishes).
The bloom is a rotation in $\delta$ at fixed scale.
Shih~\cite{Shih:2008} showed that in O'Raifeartaigh models with
fields of R-charge $\neq 0, 2$, the one-loop Coleman--Weinberg
potential gives $m_X^2 < 0$, destabilising the pseudo-modulus
away from the origin---precisely the mechanism needed to push the
seed toward the physical spectrum.


\subsection{The instanton potential and $\mathbb{Z}_6$}

The anomalous $\UU(1)_A$ of $\SU(\Nc)$ SQCD with $\Nf$ flavours
has anomaly coefficient $2\Nf$~\cite{Seiberg:1994bz,Tong:SUSY}.
Instantons break $\UU(1)_A$ to $\mathbb{Z}_{2\Nf}$;
for $\Nf = 3$ the surviving discrete symmetry is $\mathbb{Z}_6$.

In the charge parametrisation, every physical observable has the
$\mathbb{Z}_3$ permutation symmetry
$\delta \to \delta - 2\pi/3$.
The effective potential is therefore a Fourier series in
$\cos(3n\delta)$; the leading instanton term is
\begin{equation}
  V_{\rm inst} \propto -\cos(3\delta)\,,
  \label{eq:cos3delta}
\end{equation}
with three minima at $\delta = 0, 2\pi/3, 4\pi/3$.
Note that $\delta$ parametrises mass \emph{ratios} (the Cartan
direction) and is independent of the $\UU(1)_A$ phase of
$\det M$; the $\mathbb{Z}_3$ permutation symmetry constraining
$V(\delta)$ is a geometric property of the charge lattice.


\subsection{The $\zch_b$ relation}

The following empirical relation holds to remarkable precision:
\begin{equation}
  \boxed{\zch_b = 3\zch_s + \zch_c\,,}
  \label{eq:v0_doubling_intro}
\end{equation}
with $3\zch_s + \zch_c = 64.67$ vs.\ $\zch_b = 64.68$
($0.01\%$).

A \emph{possible} dynamical origin comes from bion physics.
On $\mathbb{R}^3 \times S^1$, $\SU(\Nc)$ instantons
fractionalise into $\Nc$
monopole-instantons~\cite{Unsal:2008}.
Monopole--anti-monopole (bion) cross-terms contribute to the
K\"ahler potential:
\begin{equation}
  K_{\rm bion} = \frac{\zeta^2}{\Lambda^2}
  \Bigl|\sum_k s_k \zch_k\Bigr|^2\,,
  \label{eq:Kbion}
\end{equation}
where $s_k = \pm 1$ are topological signs.

\paragraph{Caveat.}
The \"Unsal bion analysis~\cite{Unsal:2008} was developed for
QCD(adj) (adjoint fermions on $\mathbb{R}^3 \times S^1$), not
for SQCD with fundamental matter.
Equation~\eqref{eq:v0_doubling_intro} should be regarded as an
\emph{empirical relation} ($0.01\%$) with a suggestive but not
yet rigorous bion motivation.
The predicted $m_b = \zch_b^2 = 4183\MeV$
(PDG: $4183 \pm 30\MeV$).

\paragraph{Bloom dynamics: frozen-sign vs.\ adiabatic.}
The adiabatic bion potential
$V_{\rm ad} = |\sum s_k \zch_k|^2$ is exactly flat on the
energy-balance manifold.
When topological signs are \emph{frozen} at the seed
configuration, $V_{\rm frozen}$ provides a restoring force
opposing the bloom.
The equilibrium of $\cos(3\delta)$ (driving) and
$V_{\rm frozen}$ (restoring) determines the physical phase.

\paragraph{Coleman--Weinberg analysis: a negative result.}
A one-loop CW computation on the constraint surface finds that
\emph{all} perturbative contributions minimise at
$\delta = 120^\circ$ (the ``democratic'' point), not at
$\delta_{\rm phys} = 132.73^\circ$.
The physical deviation $\alpha_{\rm lep} = 12.73^\circ$ is a
\emph{flat direction} of the perturbative effective potential:
the tree-level quartic gives
$V_{\rm quartic} \propto \mathrm{Var}(1/\zch_i^2)$, and the CW
correction has $\mathrm{STr}\,\mathcal{M}^4 \propto V_{\rm quartic}$
(same angular profile).
The Cartan angle requires non-perturbative K\"ahler dynamics.


\subsection{$F$-terms are lepton masses}

When electroweak symmetry breaks, the $F$-term system of the
confined theory becomes overconstrained: 18~real equations for
12~variables.
SUSY necessarily breaks, with:
\begin{equation}
  |F_{2j}| = m_j\,,\qquad j = e, \mu, \tau\,.
  \label{eq:Fterms_leptons}
\end{equation}
The goldstino is predominantly the $\tau$-mesino (largest
$F$-term).
The gravitino mass
$m_{3/2} = |F|/(\sqrt{3}\,M_{\rm Pl})
\approx 4 \times 10^{-16}\GeV$
is negligibly small.

\paragraph{The scalar--fermion splitting.}
The universal Volkov--Akulov splitting is
$m^2_{\rm scalar} = m^2_{\rm fermion} + f_\pi^2$.
For the pion--muon pair:
$m_\pi^2 = m_\mu^2 + f_\pi^2 = 11164 + 8501 = 19665\MeV^2$,
giving $m_\pi = 140.2\MeV$ (PDG: $139.6\MeV$, $0.4\%$).


\subsection{Connection to the Seiberg dynamics}

The O'Raifeartaigh model is not an independent postulate but the
SUSY-breaking \emph{deformation} of the Seiberg vacuum.
The three constraining mechanisms work in concert:
\begin{enumerate}
  \item \textbf{Seiberg constraint} $\det M = \Lambda^6$:
    fixes the product $a_i \cdot d_i = \mu\nu$
    (family-independent).
  \item \textbf{O'Raifeartaigh seed} $gv/m = \sqrt{3}$:
    fixes the charge ratio within a seed triple
    to the algebraic value $r = 2 - \sqrt{3}$.
  \item \textbf{Instanton + bion bloom}: the $\cos(3\delta)$
    potential redistributes charge from the seed to the physical
    spectrum, while the bion restoring force constrains one mass
    combination ($\zch_b = 3\zch_s + \zch_c$).
\end{enumerate}


%======================================================================
\section{The inverse sector and chain}
\label{sec:inverse}
%======================================================================

\subsection{The dual energy-balance on $(d, s, b)$}

Evaluating on inverse masses ($\zch_j \to 1/\zch_j$):
\begin{equation}
  K^{-1}(d,s,b) = \frac{1/\zch_d^2 + 1/\zch_s^2 + 1/\zch_b^2}
  {(1/\zch_d + 1/\zch_s + 1/\zch_b)^2}
  = 0.6652\,,
\end{equation}
$0.22\%$ from $2/3$, using PDG 2024 $\overline{\rm MS}$
masses~\cite{PDG:2024}.
Under RG running, the inverse down-type triple improves
monotonically toward the UV:
\begin{center}
\begin{tabular}{lll}
\hline
$Q$ & $K^{-1}(d,s,b)$ & deviation \\
\hline
$M_Z$ & $0.66750$ & $+0.126\%$ \\
$1\TeV$ & $0.66719$ & $+0.078\%$ \\
$10\TeV$ & $0.66695$ & $+0.042\%$ \\
$100\TeV$ & $0.66675$ & $+0.012\%$ \\
$\mathbf{280\TeV}$ & $\mathbf{0.66667}$ & $\mathbf{0.000\%}$ \\
$10^6\GeV$ & $0.66657$ & $-0.014\%$ \\
$M_{\rm Planck}$ & $0.66539$ & $-0.191\%$ \\
\hline
\end{tabular}
\end{center}
The crossing $K^{-1}(d,s,b) = 2/3$ occurs at
$Q = 280\TeV$ (5-loop SMDR, interpolated).
A second inverse triple $K^{-1}(s,c,-t)$ crosses $2/3$
at $Q = 82\TeV$.
Both share the strange quark as pivot, forming an
\textbf{inverse chain}: $(d,s,b)$--$(s,c,-t)$.


\subsection{The interleaving theorem}

Given a single set of three charges $\zch_1 < \zch_2 < \zch_3$,
the Seiberg dynamics produces the \textbf{direct branch}
$a_i = \mu\zch_i^2$ and the \textbf{inverse branch}
$d_i = \nu/\zch_i^2$, with opposite orderings.

\begin{theorem}[Interleaving]
\label{thm:interleave}
The six masses form the alternating ladder
\begin{equation}
  d_1 > a_3 > d_2 > a_2 > d_3 > a_1
\end{equation}
if and only if
$\max(\zch_1^2\zch_3^2, \zch_2^4) < \nu/\mu
< \zch_2^2\zch_3^2$.
\end{theorem}

\begin{proof}
Each inequality constrains $\nu/\mu$.
$d_1 > a_3$ requires $\nu/\zch_1^2 > \mu\zch_3^2$,
i.e.\ $\nu/\mu > \zch_1^2\zch_3^2$.
$a_3 > d_2$ requires $\mu\zch_3^2 > \nu/\zch_2^2$,
i.e.\ $\nu/\mu < \zch_2^2\zch_3^2$.
The remaining four inequalities are weaker or equivalent.
The tightest lower bound is
$\max(\zch_1^2\zch_3^2, \zch_2^4)$; the upper is
$\zch_2^2\zch_3^2$.
\end{proof}

The product $a_i \cdot d_i = \mu\nu$ is family-independent
(from $\det M = \Lambda^6$)---the geometric mean constraint:
\begin{equation}
  a_i \cdot d_i = \mu\nu = \text{const}\,,
  \qquad\forall\; i\,.
  \label{eq:productrule}
\end{equation}

\paragraph{Quantitative status.}
The interleaving is \emph{proven} as abstract mathematics.
However, identifying the six physical quark masses with the six
rungs faces a quantitative obstruction: if we take down-type
charges and compute $d_2 = \nu/\zch_s^2$, we get $8638\MeV$
versus $m_c = 1273\MeV$---a factor $6.8$ discrepancy.
\textbf{The waterfall cannot reproduce the up-type spectrum from
down-type charges alone without K\"ahler corrections.}


\subsection{The chain structure}

The quark mass spectrum exhibits overlapping triples:
\begin{equation}
  \underbrace{(t,b,c)}_{\text{Fermi}} \;-\;
  \underbrace{(b,c,-s)}_{\text{transition}} \;-\;
  \underbrace{(c,s,0)}_{\text{seed}} \;-\;
  \underbrace{(s,0,m_u{+}m_d)}_{\text{lightest}}\,.
  \label{eq:chain}
\end{equation}

\begin{center}
\begin{tabular}{llll}
\hline
Link & Triple & $K$ & Status \\
\hline
1 & $(+t, +b, +c)$ & $0.6693$ ($0.40\%$) & Fermi side \\
2 & $(+b, +c, -s)$ & $0.6747$ ($1.21\%$) & sign flip on $s$ \\
3 & $(+c, +s, 0)$ & $0.6646$ ($0.31\%$) & seed triple \\
4 & $(+s, 0, \zch_{u+d})$ & $0.6649$ ($0.28\%$) & lightest seed \\
\hline
\end{tabular}
\end{center}

\paragraph{The repeated seed map.}
The chiral tail is two iterations of the seed map
$\zch_{k-1} = r\,\zch_k$ with $r = 2 - \sqrt{3}$:
$\zch_s/\zch_c = 0.271$ and $\zch_{u+d}/\zch_s = 0.271$,
both within $1.1\%$ of $r \approx 0.268$.

\paragraph{Threshold decoupling and the chiral tail.}
Below the charm threshold, heavy flavours have been integrated out
in the sequence $6 \to 5 \to 4 \to 3$ active quarks.
In this regime the natural variables are those of chiral
perturbation theory: the GOR combination
$m_u + m_d$ replaces the individual up and down masses.
The chain's tail is an EFT reorganisation of variables after
threshold decoupling.
This also gives a structural reading of the coefficient 3 in the
bion relation $\zch_b = \zch_c + 3\zch_s$: three threshold
crossings ($t \to b \to c \to s$) separate the $b$-quark from
the strange-pivoted chiral sector.

\paragraph{Top completion.}
The top link can be \emph{derived} rather than fitted.
Solving $K[\zch_t, \zch_b, \zch_c] = 2/3$ for positive $\zch_t$
is a quadratic equation whose positive root gives:
\begin{equation}
  \zch_t^{\rm pred} = 2(\zch_b + \zch_c)
  + \sqrt{3}\,\sqrt{\zch_b^2 + 4\,\zch_b\zch_c + \zch_c^2}\,.
  \label{eq:top_completion}
\end{equation}
Numerically: $\zch_t^{\rm pred} = 410.5\MeV^{1/2}$, predicting
$m_t^{\rm pred} = \zch_t^2 = 168.5\GeV$, compared with
$m_t^{\rm pole} = 172.4\GeV$ ($2.3\%$).

\paragraph{Chain compression.}
The chain is compressed by three relations:
\begin{align}
  \zch_s &\approx (2-\sqrt{3})\,\zch_c
    &&\text{(seed map)}\,,\nonumber\\
  \zch_{u+d} &\approx (2-\sqrt{3})\,\zch_s
    &&\text{(iterated seed)}\,,\nonumber\\
  \zch_b &\approx \zch_c + 3\zch_s
    &&\text{(bion / threshold)}\,,
  \label{eq:chain_compression}
\end{align}
plus the positive completion~\eqref{eq:top_completion} for
$\zch_t$.
Given $\zch_c$ as the single input, the remaining five charges
are determined to within $\sim\!2\%$.


\subsection{The cross-sector product rule and the Cabibbo angle}

If leptons and down quarks share the same three family charges
on opposite branches, the product rule $a_i \cdot d_i = \mu\nu$
gives a cross-sector prediction at the holomorphic level:
\begin{equation}
  m_e \cdot m_b^{\rm hol} = m_\mu \cdot m_s^{\rm hol}
  = m_\tau \cdot m_d^{\rm hol}\,.
  \label{eq:cross_sector}
\end{equation}
This implies $m_d^{\rm hol}/m_s^{\rm hol} = m_\mu/m_\tau$,
and hence a holomorphic prediction for the Cabibbo angle via the
Gatto--Sartori--Tonin relation~\cite{GattoSartoriTonin:1968}:
\begin{equation}
  |V_{us}|_{\rm hol} = \sqrt{m_\mu/m_\tau} = 0.244\,,
  \label{eq:Vus_hol}
\end{equation}
compared with the observed $|V_{us}| = 0.224$ ($8.7\%$ off).
However, the Cabibbo angle identity $\theta_C = \delta$
(Appendix~\ref{app:CKM}) gives the sharper prediction
$|V_{us}| = 2/9 = 0.222$ ($0.9\%$).
Combined with the Gatto--Sartori--Tonin relation, this yields
the cross-sector quark mass ratio
\begin{equation}
  \frac{m_d}{m_s} = \delta^2 = \frac{4}{81} = 0.0494\,,
  \label{eq:md_over_ms}
\end{equation}
matching PDG $0.050$ to~$1.2\%$ with zero free parameters.
Together with the inverse Koide condition $K^{-1}(d,s,b) = 2/3$
and the bion relation $\sqrt{m_b} = 3\sqrt{m_s} + \sqrt{m_c}$,
the three down-type quark masses are determined from a single input
($m_c$), giving $m_s = 91.6\MeV$ ($1.9\%$),
$m_b = 4141\MeV$ ($1.0\%$).


\subsection{The CKM matrix: an open problem}

The remaining CKM elements ($|V_{cb}|$, $|V_{ub}|$,
$\delta_{\rm CP}$) are \textbf{not predicted}.
Empirical patterns have been noted: the charge ratio
$R = \zch_b^2/\zch_c^2 = (5+\sqrt{21})/2$ satisfies the
self-dual equation $R^2 - 5R + 1 = 0$, and
$R^3 = 55 + 12\sqrt{21}$ is the fundamental unit of the real
quadratic field $\mathbb{Q}(\sqrt{21})$.
But these are \emph{observations about the data}, not derivations
from the theory.

The CP phase satisfies a suggestive $N = 3$ consistency condition:
$\cos^2(2\delta) = N/(N+4) = N/(2N+1)$ holds only for $N = 3$
(since $N + 4 = 2N + 1$ gives $N = 3$).


%======================================================================
\section{Generation uniqueness and $\SO(32)$}
\label{sec:SO32}
%======================================================================

\subsection{Why five quarks in the mesons}

The SM has six quarks but the s-confining $\SU(3)_F$ requires
$\Nf = \Nc = 3$.
The top quark ($m_t = 173\GeV \gg \Lambda_F$) is integrated out.
The remaining five quarks ($u,d,s,c,b$) decouple sequentially
(Table~\ref{tab:phases}).
The Diophantine system selects $(r, s) = (3, 2)$: three down-type
and two up-type light quarks, with $N = rs/2 = 3$ generations.

\subsection{The Diophantine system}

The sBootstrap counting~\cite{Rivero:2024epjc} imposes three
simultaneous conditions:
\begin{align}
  rs &= 2N && \text{(total flavour pairs)}\,,\nonumber\\
  \frac{r(r+1)}{2} &= 2N && \text{(symmetric composites)}\,,
  \nonumber\\
  r^2 + s^2 - 1 &= 4N && \text{(adjoint content)}\,.
\end{align}
Together they have the unique solution $(N, r, s) = (3, 3, 2)$.
The composites fill $\mathbf{24} \oplus \mathbf{15} \oplus
\overline{\mathbf{15}}$ of $\SU(5)_f$.

\subsection{Five additional $N = 3$ uniqueness criteria}

Beyond the Diophantine system, four additional criteria each
select $N = 3$:

\begin{enumerate}
  \item \textbf{Discriminant matching}:
    the discriminant of the characteristic polynomial of $\SU(N)$
    gives $D_{\rm alg} = (N+2)^2 - 4$, while the dimension of
    the s-confining constraint surface is $D_{\rm dyn} = N(2N+1)
    = \dim(\Sp(2N))$.
    Requiring $D_{\rm alg} = D_{\rm dyn}$ gives
    $(N+2)^2 - 4 = N(2N+1)$, which has the unique solution
    $N = 3$ ($21 = 21$).

  \item \textbf{Binomial identity}:
    $\binom{3N}{N} = (N+1) \cdot N(2N+1)$ holds only at $N = 3$:
    $84 = 4 \times 21$.
    The number $84 = \binom{9}{3}$ counts the $C_3$ polarisations
    of 11d SUGRA under $\SO(9)$, and 21 is
    $\dim(\SU(3)_F \text{ composites at } \Nf = \Nc = 3)$.

  \item \textbf{CP-phase consistency} (conditional):
    requiring $\cos^2(2\delta) = N/(N+4)$ and
    $\cos^2(2\delta) = N/(2N+1)$ simultaneously gives
    $N + 4 = 2N + 1$, which requires $N = 3$.
    This uses the observed CP phase, not predicted by the
    framework; it is a consistency check.

  \item \textbf{Casimir--Koide identity}:
    for $\UU(N)$ with canonical normalisation, the energy-balance
    ratio generalises to $K_N = 2/N$ (Theorem~\ref{thm:n2}
    generalised to $N$ charges).
    The mean-square weight $\langle J_3^2\rangle
    = N(N^2-1)/(12N) = (N^2-1)/12$ of the principal
    $\mathfrak{sl}(2)$ embedding in
    $\mathfrak{sl}(N)$~\cite{Kostant:1959}
    (where $J_3$ is the weight operator in the fundamental
    representation).
    The condition $K_N = \langle J_3^2\rangle$ gives
    $2/N = (N^2-1)/12$, i.e.\ $N(N^2-1) = 24$.
    Factoring: $(N-3)(N^2+3N+8) = 0$; the quadratic has
    discriminant $\Delta = 9 - 32 < 0$, so the unique positive
    integer solution is $N = 3$.
    Equivalently, at $N = 3$ and only $N = 3$, the representation
    dimension $2j + 1 = N$ coincides with the algebra dimension
    $\dim\,\mathfrak{sl}(2) = 3$, giving
    $C_2/N = C_2/3 = \langle J_3^2\rangle = K = 2/3$.
\end{enumerate}

\subsection{$\SO(32)$ selection}

The ten-dimensional heterotic string has two anomaly-free gauge
groups: $\SO(32)$ and $E_8 \times E_8$.
The sBootstrap composite content selects between them.

$\SO(32) \supset \SU(16) \times \UU(1) \supset \SU(5) \times
\SU(3) \times \UU(1)^2$.
The adjoint $\mathbf{496}$ contains $\mathbf{24}$ and
$\mathbf{15} + \overline{\mathbf{15}}$---precisely the
sBootstrap content.
The $Q = 4/3$ states are projected out by the $\SU(3)_L$
decomposition.

Every decomposition of $E_8$ under $\SU(5)$ produces only the
antisymmetric $\mathbf{10}$, never the symmetric
$\mathbf{15}$~\cite{Rivero:2024epjc}.
The sBootstrap therefore selects $\SO(32)$ over
$E_8 \times E_8$.

\begin{table}[ht]
\centering\small
\caption{$\SU(5)_f$ composite content and sBootstrap scorecard.}
\label{tab:SU5_scorecard}
\begin{tabular}{llll}
\hline
$\SU(5)_f$ rep & SM decomposition & Role & Status \\
\hline
$\mathbf{24}$ & $(\mathbf{8},\mathbf{1})_0 + \ldots$ & gauge + Higgs & \\
  & $+ (\mathbf{3},\mathbf{2})_{-5/6} + \text{c.c.}$ & $Q{=}4/3$ proj.\ & \\
$\mathbf{15}$ & $(\mathbf{1},\mathbf{3})_{-1}$ & sleptons & $\checkmark$ \\
  & $(\mathbf{3},\mathbf{2})_{1/6}$ & leptoquarks & $\approx$ \\
  & $(\mathbf{6},\mathbf{1})_{-2/3}$ & diquarks & $\times$ \\
\hline
\multicolumn{4}{l}{84 DOF: exact 16 (19\%), approx.\ 44 (52\%), obstr.\ 24 (29\%).} \\
\hline
\end{tabular}
\end{table}

In 11d SUGRA, $\mathbf{6}_C$ belongs to
$\mathrm{Sym}^2(9) = \mathbf{45} \subset g_{MN}$ (gravity
sector), not $\Lambda^3(9) = \mathbf{84} \subset C_3$ (gauge
sector), explaining the colour-sextet obstruction.


%======================================================================
\section{The effective Lagrangian}
\label{sec:lagrangian}
%======================================================================

\subsection{Scenario (a): sequential integration}

\paragraph{The UV theory.}
The starting point is $\SU(3)_F$ SQCD with $\Nf = 5$ massive
flavours $(u,d,s,c,b)$.
As the RG scale decreases below $m_b$ and $m_c$, heavy flavours
are integrated out sequentially (Table~\ref{tab:phases}).
At each threshold the dynamical scale matches:
$\Lambda_3^6 = \Lambda_4^5 m_c$,
$\Lambda_4^5 = \Lambda_5^4 m_b$.

\paragraph{The $\Nf = 4$ ISS window.}
Between $m_c$ and $m_b$, the theory has $\Nf = 4$ and sits in the
ISS window~\cite{ISS:2006}.
The IR-free magnetic dual has gauge group
$\SU(\Nf - \Nc) = \SU(1) = \text{trivial}$, with the macroscopic
superpotential:
\begin{equation}
  W_{\rm ISS} = h\,\Tr\,\varphi\,\Phi\,\tilde\varphi
  - h\mu^2\,\Tr\Phi\,,
  \label{eq:W_ISS}
\end{equation}
where $\Phi_{ij}$ is the meson, $\varphi$ and
$\tilde\varphi$ are magnetic quarks, $h$ is the magnetic Yukawa,
and $\mu^2 \sim m\Lambda$ is the SUSY-breaking parameter.
SUSY is broken by the rank condition.
At $\Nf = 3$ the magnetic quarks disappear and
eq.~\eqref{eq:W_ISS} reduces to the s-confined
superpotential---the ISS window serves only as a \emph{transient}
SUSY-breaking phase during the RG descent.

\paragraph{The $\Nf = 3$ confined theory.}
Below $m_c$, the theory s-confines with the effective
superpotential:
\begin{equation}
  W_{\rm conf} = X\bigl(\det M^{(3)} - B\bar{B} - \Lambda_3^6\bigr)
    + \Tr(\hat{m}\,M^{(3)})\,,
  \label{eq:W_conf}
\end{equation}
where $M^{(3)}$ is the $3 \times 3$ light-flavour meson matrix,
$X$ is the Lagrange multiplier, and
$\hat{m} = \diag(m_u, m_d, m_s)$ are the light preon masses.

\paragraph{The full superpotential.}
Combining all sectors:
\begin{align}
  W_{\rm (a)} &= X\bigl(\det M^{(3)} - B\bar{B} - \Lambda_3^6\bigr)
    + \Tr(\hat{m}\,M^{(3)})
    \nonumber\\
  &\quad + c_3\,\frac{(\det M^{(3)})^3}{\Lambda_3^{18}}
    \quad\text{(three-instanton)}
    \nonumber\\
  &\quad + y_c\,H_u\,M^c{}_c + y_b\,H_d\,M^b{}_b
    \quad\text{(heavy-flavour Yukawas)}
    \nonumber\\
  &\quad + y_t\,H_u\,Q^t\bar{Q}_t
    \quad\text{(top: partial compositeness)}\,.
  \label{eq:W_scenario_a}
\end{align}

\paragraph{The K\"ahler potential.}
\begin{equation}
  K_{\rm (a)} = \Tr(M^\dagger M) + |X|^2 + |B|^2 + |\bar{B}|^2
  - \frac{|X|^4}{12\mu^2}
  + \frac{\zeta^2}{\Lambda^2}
    \Bigl|\sum_k s_k \zch_k\Bigr|^2\,.
  \label{eq:Kahler_a}
\end{equation}
The $-|X|^4/(12\mu^2)$ term stabilises the pseudo-modulus at
$gv/m = \sqrt{3}$.
The soft potential $V_{\rm soft} = f_\pi^2\,\Tr(M^\dagger M)$
gives universal splitting
$m^2_{\rm scalar} - m^2_{\rm fermion} = f_\pi^2$.

\paragraph{Parameter count.}
Three UV preon masses (or equivalently $\mu$, $\alpha$, and the
product $P$); the dynamical scale $\Lambda_F$; the adjoint mass
$\mu_\Phi$; and the K\"ahler corrections $Z_i$.
Of these, $\Lambda_F$ is fixed by matching to $\alpha_s(M_Z)$;
$\mu_\Phi$ is empirical ($\lesssim 4.5\TeV$); and the $Z_i$
are the master unknowns.


\subsection{Scenario (b): speculative Pati--Salam extension}

Scenario~(b) requires a Pati--Salam gauge group ($\Nc = 4$,
$\Nf = 6$) that is NOT part of the core sBootstrap framework.
In this scenario:
\begin{enumerate}
  \item The $\mathcal{N}{=}2$ $\SU(\Nc)$ theory with $\Nf$
    hypermultiplets has superpotential~\cite{HMS:1997,Fayet:1976}
    $W = \sqrt{2}\,g\,Q^i\Phi\tilde{Q}_i + \sqrt{2}\,m^j_i
    Q^i\tilde{Q}_j$, where the Yukawa $\sqrt{2}\,g$ equals the
    gauge coupling by $\mathcal{N}{=}2$ SUSY.
  \item Adding an adjoint mass
    $\Delta W = \mu_\Phi\,\Tr\Phi^2$ breaks
    $\mathcal{N}{=}2 \to \mathcal{N}{=}1$ and generates the
    quartic~\eqref{eq:quartic}.
  \item At the baryonic root~\cite{Csaki:2011xn}, the magnetic
    adjoint VEV breaks the gauge group to
    $\SU(\Nf - \Nc) = \SU(2)$, which IS $\SU(2)_L$.
\end{enumerate}
The strength is that the Yukawa hierarchy is a hierarchy of VEVs
(Seiberg seesaw), not couplings.
The limitation: the Pati--Salam embedding requires gauge structure
beyond $\SU(3)_F$, and the K\"ahler metric is still needed.
By Fradkin--Shenker complementarity, the ``UV $\mathcal{N}{=}2$''
picture and the ``emergent $\mathcal{N}{=}2$''
picture may describe the same physics, smoothly connected.


%======================================================================
\section{Discussion}
\label{sec:discussion}
%======================================================================

\paragraph{M-theory lift.}
In the Hanany--Witten brane
construction~\cite{HananyWitten:1997}, $\Nc$ D4-branes are
suspended between two NS5-branes, with $\Nf$ D6-branes providing
quarks.
Mesons are open D6--D6 strings; baryons are wrapped D4-branes.
Lifting to M-theory: D4-branes become M5-branes, D6-branes
become Taub-NUT centres.
The meson lifts to an M2-brane coupling to $C_3$ (Dirac mass).
The baryon lifts to an M5-brane coupling to $C_6$ (Majorana mass).
The involution $M_i \to \Lambda^4/M_i$ exchanges M2 and M5
branes: electromagnetic duality $C_3 \leftrightarrow C_6$.

\begin{table}[ht]
\centering\small
\caption{M-theory dictionary for $\SU(3)$ SQCD composites.}
\label{tab:Mtheory}
\begin{tabular}{lllll}
\hline
Field theory & IIA & M-theory & $C$ & Mass \\
\hline
Meson $Q\tilde{Q}$ & D6--D6 & M2 & $C_3$ & Dirac \\
Baryon $\varepsilon Q^3$ & D4 & M5 & $C_6$ & Majorana \\
$\adj(M)$ & thr.\ D4 & thr.\ M5 & mixed & branch \\
Involution & T-dual & EM dual & $C_3{\leftrightarrow}C_6$ & --- \\
\hline
\end{tabular}
\end{table}

The membrane dimensionality provides a consistency check: for a
$p$-brane with $m_i \propto \zch_i^p$,
$K_p$ is direction-independent only for $p = 2$ (the M2-brane),
by Theorem~\ref{thm:n2}.

\paragraph{The spin-2 connection.}
The $Q_2$ operator $R$ is the spin-2 component of $\mathrm{End}(V)$
under the principal $\mathfrak{sl}(2)$
(Sec.~\ref{sec:N2}).
In 11d SUGRA: $C_3 \in \Lambda^3(9) = \mathbf{84}$ carries
spin-1 (meson/M2 sector), while
$g_{MN} \in \mathrm{Sym}^2(9) = \mathbf{45}$ carries spin-2
(gravitational sector).
Colour sextets live in
$\mathrm{Sym}^2(\mathbf{3}_C) = \mathbf{6}_C \subset \mathbf{45}$,
in the spin-2/graviton sector.
The $Q_2$ rotation $R$, being spin-2, may therefore be a
\emph{gravitational} operation---a metric fluctuation connecting
the $C_3$ (meson/lepton) sector to the $g_{MN}$ (quark) sector.
This is speculative but structurally consistent: the spin quantum
number of $R$ matches the spin of the graviton.

The number $84 = \binom{9}{3}$ appears in five contexts:
(i)~$C_3$ polarisations ($\Lambda^3(\mathbb{R}^9)$);
(ii)~$E_8 \to \SU(9)$: $248 = 80 + 84 + \overline{84}$;
(iii)~SM fermions: $96 = 84 + 12$;
(iv)~number theory: $|\text{disc}(\mathbb{Q}(\sqrt{-21}))| = 84$;
(v)~the binomial identity $\binom{9}{3} = 4 \times 21$.

\paragraph{Composite $\SU(2)_L$ (speculative).}
The CST construction~\cite{Csaki:2011xn} offers a mechanism in
which $\SU(2)_L$ is a residual magnetic gauge group at the
baryonic root, but requires a Pati--Salam embedding
($\Nc = 4$, $\Nf = 6$) beyond the core $\SU(3)_F$ framework.
In the core sBootstrap, $\SU(2)_L$ is an
\textbf{elementary} gauge symmetry in $G_{\rm UV}$.

\paragraph{Holomorphic vs.\ physical masses.}
The meson masses from the Seiberg vacuum are holomorphic,
protected by holomorphy.
Physical masses are $m_{\rm phys} = m_{\rm hol}/\sqrt{Z_i}$.
For leptons, $Z$ is generation-universal (to $0.001\%$).
For quarks, $m_c/m_s$ does NOT run ($\gamma_m$ is
flavour-independent to 5-loop), and the lattice value
$m_c/m_s = 11.783 \pm 0.025$ differs from the lepton-derived
prediction $(2+\sqrt{3})^2 = 13.93$ by $18\%$.
This discrepancy dissolves when sectors are allowed different
Cartan angles: the chain seed gives
$\zch_s/\zch_c = 0.271$, within $1.1\%$ of $r = 2 - \sqrt{3}$.

\paragraph{Hierarchy reframing.}
In the sBootstrap, SUSY is broken at $f = f_\pi \sim 92\MeV$
---the QCD scale, not the TeV scale.
The electroweak VEV $v_{\rm EW} = 246\GeV$ is derived:
$v_{\rm EW} \sim m_t / y_t \sim 173/0.7 \sim 247\GeV$.
The apparent hierarchy $v_{\rm EW}/\Lambda_{\rm QCD} \sim 10^3$
is the Yukawa hierarchy---the flavour problem that the
charge-squared law addresses.
This is a \emph{reframing}, not a solution.

\paragraph{Strong CP and Nelson--Barr.}
Identifying the involution with CP gives a
Nelson--Barr~\cite{Nelson:1984,Barr:1984} solution.
At tree level: $\arg\det m_q = 0$.
Radiative corrections:
$\delta\bar\theta \sim (\alpha/4\pi)^2 J \sim 10^{-11}$,
nine orders of magnitude below $|\bar\theta| < 10^{-10}$.

\paragraph{The AM-GM scale chain.}
Three mass scales from $\SU(3)$ confinement:
\begin{align}
  M_0 &= \mathrm{AM}^2 = 313.8\MeV
    && (\text{Koide scale} \approx m_p/3)\,,\\
  \tfrac{1}{2}f_\pi &= \mathrm{GM}^2 = 45.8\MeV
    && (\text{SUSY breaking})\,,\\
  \mu &= 6\,M_0 = 1883\MeV
    && (\text{mass scale} \approx 2\,m_p)\,.
\end{align}
The $f_\pi$ relation $f_\pi = 2(\zch_e\,\zch_\mu\,\zch_\tau)^{2/3}$
holds to $0.7\%$.
The sum $\mu = m_e + m_\mu + m_\tau = 1883\MeV$ coincides with
$2\,m_p = 1877\MeV$ ($0.3\%$).
The probability of some match at $0.3\%$ among $\sim\!50$ hadron
masses is $p \approx 15\%$ after look-elsewhere correction.

\paragraph{Adjugate crossing.}
Under RG running, $K(e,\mu,\tau)$ and $K^{-1}(d,s,b)$ cross at
$Q \approx 15.7\GeV$: $K_{\rm lep} = K^{-1}_{\rm down} = 0.6663$,
and the gap at $m_b$ is only $0.06\%$.
This crossing scale is near the $b$-quark threshold, consistent
with the $\Nf = 4 \to \Nf = 3$ phase transition.


%======================================================================
\section{Open problems and conclusions}
\label{sec:open}
%======================================================================

\subsection{Numerical results}

\begin{table}[ht]
\centering\small
\begin{tabular}{lllll}
\hline
Observable & Method & Predicted & PDG 2024 & Status \\
\hline
$K(e,\mu,\tau)$ & trace id.\ & $2/3$ & $0.666661$ & $0.001\%$ \\
$K(-\pi,D_s,B)$ & meson & $\approx 2/3$ & $0.6674$ & $0.10\%$ \\
$K(-\pi,D,B)$ & meson & $\approx 2/3$ & $0.673$ & $0.91\%$ \\
$K^{-1}(d,s,b)$ & seesaw & $\approx 2/3$ & $0.6652$ & $0.22\%$ \\
$m_b$ (bion) & eq.~\eqref{eq:v0_doubling_intro} & $4183\MeV$ & $4183$ & $0.01\%$ \\
$m_t$ (compl.) & eq.~\eqref{eq:top_completion} & $168.5\GeV$ & $172.4$ & $2.3\%$ \\
$|V_{us}|$ & $\sqrt{m_\mu/m_\tau}$ & $0.244$ & $0.224$ & $8.7\%$ \\
\hline
\end{tabular}
\caption{Quantitative results.
  All entries use PDG 2024~\cite{PDG:2024} inputs.}
\label{tab:results}
\end{table}

\paragraph{Tier~1: high precision ($< 0.1\%$).}
$K(e,\mu,\tau) = 0.666661$ ($0.001\%$) stands alone as the
most precise mass relation in particle physics outside the
gauge sector.
The bion prediction $\zch_b = 3\zch_s + \zch_c$ ($0.01\%$) is
also in this tier.
Joint $p < 10^{-7}$.

\paragraph{Tier~2: moderate precision ($0.1$--$1\%$).}
The meson triple $K(-\pi, D_s, B)$ ($0.10\%$) and the
inverse triple $K^{-1}(d,s,b)$ ($0.22\%$).
At $Q = 100\TeV$, $K^{-1}(s,c,-t) = 0.6666$ ($0.004\%$) joins
Tier~1.

\paragraph{Tier~3: the chain as a correlated set.}
The quark chain is individually $0.3$--$1.2\%$ per link, but
the probability that random masses produce four consecutive
windows all within $1.5\%$ of $2/3$ is $p \sim 1/21\,000$.

\paragraph{Prediction vs.\ postdiction.}
The lepton Koide (Koide 1983) and the inverse quark Koide
(Rivero 2005) are postdictions.
The framework's genuine predictions are the LFV ratio and the
family gauge boson masses.


\subsection{Open problems}
\label{sec:open_problems}

\paragraph{Resolved or partially resolved.}
\begin{itemize}
\item \textbf{The fermion K\"ahler metric} (formerly ``master
  obstruction'').
  The emergent $\mathcal{N}{=}2$ hypermultiplet structure
  protects the fermion K\"ahler metric: the $\mathcal{N}{=}2$
  breaking correction is $\delta Z/Z \sim g^2\Lambda^2/\mu_\Phi^2
  \sim 5 \times 10^{-9}$, negligible.
  This \emph{explains} the $0.001\%$ lepton Koide: it is
  $\mathcal{N}{=}2$ protection, not a mystery.
  \textbf{This part is solved.}

\item \textbf{The $m_c/m_s$ ratio} (formerly part of Problem~1).
  The ``discrepancy'' dissolves when sectors are allowed
  different Cartan angles: the chain seed gives
  $\zch_s/\zch_c = 0.271$, within $1.1\%$ of $r = 2 - \sqrt{3}$.
  No K\"ahler correction is needed.

\item \textbf{The $Q_2$ operator} (formerly ``open: can the
  $\mathbb{Z}_2$ be promoted?'').
  The operator $R$ is explicitly constructed
  (Sec.~\ref{sec:N2}).
  It satisfies $R \cdot \mathbf{M} = 2\,\adj(\mathbf{M})$
  and the cubic algebra
  $R^3 = \Tr(D^2)\,R + 2\Lambda^6\,I$.
  $R$ is the pure spin-2 component of $\mathrm{End}(V)$ under
  Kostant's principal $\mathfrak{sl}(2)$.
  The normalised operator is an ``almost-involution'' with
  eigenvalues $\approx \{+1.07,\; 0,\; -1.07\}$.
\end{itemize}

\paragraph{Remaining open problems.}
\begin{enumerate}
\item \textbf{The Cartan angle $\alpha \approx 12.7^\circ$}
  (absorbs old Problems 3 and 4).
  One-loop CW analysis shows all perturbative contributions
  minimise at $\delta = 120^\circ$, not at
  $\delta_{\rm phys} = 132.73^\circ$.
  The deviation is a flat direction of the perturbative theory.
  Its value requires the non-perturbative K\"ahler metric for
  the \emph{scalar} meson VEV (not the fermion metric, which is
  solved).

\item \textbf{The CKM matrix.}
  The off-diagonal meson mass matrix block-decouples from the
  diagonal mesons at tree level.
  CKM mixing cannot arise from tree-level meson mixing;
  it requires K\"ahler corrections, loop effects, or RG running.
  $|V_{us}|_{\rm hol} = 0.244$ is $8.7\%$ off.

\item \textbf{PMNS and neutrino masses.}
  The baryons $B, \tilde{B}$ are SM singlets.
  Tree-level PMNS is diagonal---contradicting data.

\item \textbf{The cubic $Q_2$ algebra.}
  Eq.~\eqref{eq:R_cubic} shows $R$ satisfies a cubic.
  Does this define a consistent deformed $\mathcal{N}{=}2$
  superalgebra?

\item \textbf{The $\SO(32)$ projection.}
  The $\mathbf{496}$ contains 442 states beyond the sBootstrap
  content; the projection mechanism is unknown.
\end{enumerate}


\subsection{Falsification conditions}

\begin{itemize}
  \item The LFV ratio
    $\mathrm{Br}(\tau \to e\gamma)/\mathrm{Br}(\tau \to \mu\gamma)
    = (\zch_e/\zch_\mu)^4 \approx 5.5 \times 10^{-10}$ is
    testable at Belle~II.
  \item Family gauge bosons at $10^2$--$10^3\TeV$ produce
    flavour-changing effects testable at FCC-hh.
  \item Lattice K\"ahler metric strongly generation-dependent
    for leptons ($Z_\mu/Z_e \neq 1$ at $1\%$) would kill the
    framework.
  \item Discovery of a fourth generation contradicts $N = 3$
    uniqueness.
\end{itemize}


\subsection{Conclusions}

We have presented a framework in which Standard Model fermions
are identified with composite fermions of a confining $\SU(3)_F$
family gauge theory: mesinos are leptons, and the adjugate
composites provide quark partners through partial compositeness.
The confined spectrum automatically produces a pair
$(M_i, \adj(M)_i)$ per family with the mathematical structure of
an $\mathcal{N}{=}2$ hypermultiplet---emergent from confinement,
not imposed in the UV.

The framework rests on one central ansatz---the charge-squared
law $m_i = \mu\,\zch_i^2$---which is the unique monomial
compatible with direction-independent $K = 2/3$
(Theorem~\ref{thm:n2}), supported by the Coulomb self-energy
of the composite~\cite{Koide:1983qe} and by the bilinear
structure $M = Q\tilde{Q}$.
We have been honest that this is the best-motivated ansatz, not
a theorem of the confining theory.

What is \textbf{established}:
\begin{enumerate}
  \item \textbf{Composites:} 9 mesinos furnish $3L_i + 3e_i^c$
    (exact group theory).
    Both Higgs doublets from one meson matrix:
    $H_d = M^{1,2}$, $H_u = \adj(M)^{1,2}$, ratio fixed by the
    Seiberg constraint.

  \item \textbf{The adjugate mechanism:} $M \propto 1/m$ and
    $\adj(M) \propto m$ provide two Yukawa operators from one
    composite.
    A single UV condition produces both standard lepton and
    inverse quark $K = 2/3$.

  \item \textbf{Energy-balance observations:}
    $K(e,\mu,\tau) = 0.666661$ ($0.001\%$);
    $K(-\pi, D_s, B) = 0.6674$ ($0.10\%$);
    $K^{-1}(d,s,b) = 0.6652$ ($0.22\%$);
    and at $Q = 100\TeV$, $K^{-1}(s,c,-t) = 0.6666$ ($0.004\%$).
    The quark chain has $K$ within $1.2\%$ in every link
    (joint $p \sim 1/21\,000$), compressed by three algebraic
    relations plus a completion formula predicting
    $m_t = 168.5\GeV$ ($2.3\%$).

  \item \textbf{The $Q_2$ operator:} explicitly constructed
    as the spin-2 component of $\mathrm{End}(V)$ under the
    principal $\mathfrak{sl}(2)$~\cite{Kostant:1959}, satisfying
    the cubic algebra~\eqref{eq:R_cubic} and the democratic
    identity~\eqref{eq:DRD}.

  \item \textbf{Generation uniqueness:} five criteria select
    $N = 3$, including the Casimir--Koide identity
    $N(N^2-1) = 24$; composite content selects $\SO(32)$.
\end{enumerate}

What the framework does \textbf{not yet} provide:
the full CKM matrix beyond $|V_{us}|_{\rm hol}$;
the PMNS matrix and neutrino masses;
the origin of $\alpha \approx 12.7^\circ$;
a derivation of $m = \zch^2$ from dynamics;
the K\"ahler metric $Z_i$.

\paragraph{What ``compute the K\"ahler metric'' means.}
The K\"ahler metric $Z_i$ is the wave-function renormalisation
of the meson composite $M_i = Q^i\tilde{Q}_i$.
It enters the kinetic term $K = Z_i|M_i|^2 + \cdots$ and
relates holomorphic and physical masses via
$m_{\rm phys} = m_{\rm hol}/\sqrt{Z_i}$.
If $Z_1 = Z_2 = Z_3$ (generation-universal), then $K = 2/3$ in
holomorphic masses implies $K = 2/3$ in physical masses.
The lepton data constrains $|Z_i/Z_j - 1| < 0.002\%$.

A lattice measurement would target the two-point correlator
$\langle M_i(x)\,M_i^\dagger(0)\rangle$ at zero momentum; the
residue gives $Z_i$.
One then checks whether $Z$ depends on the preon mass $m_i$.

However, this requires \textbf{lattice SQCD}, not lattice QCD:
$\mathcal{N}{=}1$ supersymmetry is explicitly broken by the lattice
discretisation, and restoring it in the continuum limit demands
fine-tuning a scalar mass parameter.
$\SU(3)$ with $\Nf = 3$ fundamentals has never been simulated
on a SUSY-preserving lattice.

The situation is actually better than this suggests.
The emergent $\mathcal{N}{=}2$ structure
(Section~\ref{sec:N2}) protects the \emph{fermion} K\"ahler
metric: $\delta Z/Z \sim g^2\Lambda^2/\mu_\Phi^2
\sim 5 \times 10^{-9}$ (Section~\ref{sec:open}).
\textbf{The fermion $Z_i$ is generation-universal; this part of
the problem is solved.}
What remains is the \emph{scalar} meson VEV correction:
the physical quark masses depend on $\langle M_i\rangle$, which
receives K\"ahler potential corrections to its scalar potential.
These corrections determine the Cartan angle $\alpha$
(Problem~1) and are the same non-perturbative dynamics that
the Coleman--Weinberg analysis (Section~\ref{sec:susy_breaking})
showed cannot be computed perturbatively.

The honest summary: a lattice computation of $Z_i$ would
\emph{confirm} generation-universality (already established by
the lepton Koide + $\mathcal{N}{=}2$ protection) but would not
resolve the Cartan angle, which requires the scalar potential
corrections at the confinement scale---a different and harder
computation.

On the experimental side, the LFV ratio is testable at Belle~II,
and family gauge bosons at $10^2$--$10^3\TeV$ produce
flavour-changing effects accessible to a future $100\TeV$ collider.


%======================================================================
\appendix
%======================================================================

\section{Seiberg--Witten curve at the Koide point}
\label{app:SW}

The $\mathcal{N}{=}2$ Seiberg--Witten curve for $\SU(3)$ with
$\Nf = 3$ massive hypermultiplets is
\begin{equation}
  y^2 = P(x)^2 - 4\Lambda^3\,Q(x)\,,
  \label{eq:SWcurve}
\end{equation}
where
\begin{equation}
  P(x) = x^3 - u_2\,x - u_3\,,\qquad
  Q(x) = \prod_{i=1}^{3}(x + m_i)\,.
\end{equation}
Here $u_2 = \frac{1}{2}\mathrm{Tr}\,\Phi^2$ and
$u_3 = \det\Phi$ are the Coulomb branch moduli of the adjoint
scalar~$\Phi$, while $m_i$ are the hypermultiplet masses.

\paragraph{Identification from the cubic $Q_2$ algebra.}
The Cayley--Hamilton equation of the traceless matrix~$\Phi$,
$\Phi^3 - u_2\,\Phi - u_3 = 0$,
coincides with the cubic algebra~(\ref{eq:R_cubic}) under
\begin{equation}
  \boxed{u_2 = \mathrm{Tr}(D^2) = \sum m_i^2\,,\qquad
  u_3 = 2\det(D) = 2\,m_1 m_2 m_3\,.}
  \label{eq:u2u3}
\end{equation}
For the lepton masses ($m_e = 0.511$, $m_\mu = 105.658$,
$m_\tau = 1776.86$~MeV):
$u_2 = 3.168\times 10^6~\mathrm{MeV}^2$,
$u_3 = 1.919\times 10^5~\mathrm{MeV}^3$.

The eigenvalues of~$\Phi$ are the roots of
$r^3 - u_2\,r - u_3 = 0$:
\begin{equation}
  \phi_1 = +1780.0\,,\quad
  \phi_2 = -0.061\,,\quad
  \phi_3 = -1780.0\;\mathrm{MeV}\,.
\end{equation}
These satisfy $\sum\phi_i = 0$ (traceless) and
$\sum\phi_i^2 = 2u_2$ (checked).
The near-vanishing of~$\phi_2 \approx 0$ signals proximity to
an Argyres--Douglas-like point where $\det\Phi \to 0$.

\paragraph{Democratic limit.}
For $m_1 = m_2 = m_3 = m$, the characteristic polynomial
factors exactly:
\begin{equation}
  P(x) = (x+m)^2(x - 2m)\,,
\end{equation}
and the sextic becomes
\begin{equation}
  f(x) = (x+m)^3\bigl[(x+m)(x-2m)^2 - 4\Lambda^3\bigr]\,.
  \label{eq:democratic_sextic}
\end{equation}
The bracket cubic has discriminant
$\Delta = 432\,\Lambda^3(m^3 - \Lambda^3)$,
which vanishes at $\Lambda = m$.
At this point a double root appears at $x = 0$:
a BPS monopole becomes massless at the democratic mass
$m = s_1/3 \approx 628$~MeV.
This is the only singularity in the democratic limit.

\paragraph{Koide constraint as codimension-3 slice.}
The curve~(\ref{eq:SWcurve}) has five parameters
$(u_2, u_3, m_1, m_2, m_3)$.
The Koide formula $K = 2/3$ imposes one constraint,
$\sum m_i = \frac{2}{3}(\sum\sqrt{m_i})^2$,
while the overall scale~$k = \sum m_i$ is fixed by data.
The remaining freedom is a single parameter: the Cartan
angle~$\delta$.
The SW curve at the Koide point is therefore a
\textbf{one-parameter family} $\mathcal{C}(\delta)$.

No topological transition, cusp, or branch-point collision
occurs at $\delta = 2/9$; the Cartan angle enters the geometry
smoothly through~$u_2(\delta)$ and~$u_3(\delta)$.
Full numerical results are in \texttt{scripts/sw\_curve\_su3\_nf3.py}
and the accompanying figures.

%----------------------------------------------------------------------
\section{The cubic $Q_2$ algebra as $\mathcal{N}{=}2$
         Cayley--Hamilton}
\label{app:cubic}

The $Q_2$ operator $R_{ij} = \langle M_k\rangle$ satisfies (Sec.~\ref{app:cubic})
\begin{equation}
  R^3 = \mathrm{Tr}(D^2)\,R + 2\Lambda^6\,I\,.
  \label{eq:cubic_Q2_app}
\end{equation}
Rearranging: $R^3 - \mathrm{Tr}(D^2)\,R - 2\det(D)\,I = 0$.
This is the Cayley--Hamilton equation for a traceless $3\times 3$
matrix with quadratic Casimir $\frac{1}{2}\mathrm{Tr}(A^2) =
\mathrm{Tr}(D^2)$ and cubic Casimir $\det(A) = 2\det(D)$.

In $\mathcal{N}{=}2$ $\SU(3)$ SQCD, the adjoint scalar~$\Phi$ is
traceless and its vacuum expectation values parameterise the
Coulomb branch. Its Cayley--Hamilton equation is precisely
Eq.~(\ref{eq:cubic_Q2_app}) with the identifications~(\ref{eq:u2u3}).

\paragraph{Physical interpretation.}
The cubic is the \emph{residual} $\mathcal{N}{=}2$ constraint
after $\mathcal{N}{=}2 \to \mathcal{N}{=}1$ breaking by the
adjoint mass~$W = \mu_\Phi\,\mathrm{Tr}(\Phi^2)$.
When $\mu_\Phi \to \infty$, the adjoint is frozen and its
eigenvalues become algebraic functions of the meson
eigenvalues~$m_i$.
The cubic changes from a dynamical equation (vacua on the Coulomb
branch) to a constraint (the meson masses must satisfy the
Cayley--Hamilton identity).

\paragraph{CKM invariance.}
Both $\mathrm{Tr}(D^2)$ and $\det(D)$ are invariant under
unitary basis changes $D \to V D V^\dagger$,
so the cubic is \textbf{CKM-invariant}:
its coefficients do not change when the family basis is rotated
by the CKM matrix.
Off-diagonal (CKM-generating) corrections therefore enter only
through \emph{deformations} of the cubic, not modifications of
its leading coefficients.

%----------------------------------------------------------------------
\section{Off-diagonal corrections and CKM}
\label{app:CKM}

At tree level the meson vacuum is flavour-diagonal
($\langle M_{ij}\rangle = 0$ for $i \neq j$),
as confirmed by Masiero and
Veneziano~\cite{MasieroVeneziano:1985} (their eq.~3.25).
CKM mixing must therefore arise from radiative corrections.

\paragraph{Masiero--Veneziano mechanism.}
In Ref.~\cite{MasieroVeneziano:1985}, the charged lepton mass
is generated by a one-loop diagram (their Fig.~2) involving a
photino propagator and the off-diagonal meson propagator
$\langle \pi_{23}\pi_{32}\rangle$:
\begin{equation}
  m_{e^-} = \frac{e^2}{12\pi^2}\,m_{\tilde\gamma}\,
  \frac{\tilde\mu_{33}}{\tilde\mu}\,
  \log\frac{m_+^2}{m_-^2}\,.
  \label{eq:MV_loop}
\end{equation}
The key point: even though $\langle M_{ij}\rangle = 0$, the
two-point function $\langle M_{ij}M_{ji}\rangle \neq 0$
(set by the scalar mass matrix).

\paragraph{Generalisation to CKM.}
In our framework, the quarks enter through partial compositeness
with Yukawa couplings $Y_i$.
The CKM-generating loop has the structure
$Y_i \times (\text{gaugino} + \text{off-diagonal meson loop})
\times Y_j$.

The leading one-loop estimate gives
\begin{equation}
  \frac{\delta M_{ij}}{M_j} \sim
  \frac{g^2}{16\pi^2}\,\frac{M_i}{\mu_\Phi}\,
  \log\frac{\mu_\Phi}{\Lambda}\,,
\end{equation}
which for $\mu_\Phi \sim 10^2$--$10^4$~MeV yields
$V_{us}^{\mathrm{1-loop}} \sim 10^{-5}$,
\textbf{short by a factor of~$\sim 10^3$}.

\paragraph{Holomorphic product rule.}
The best within-framework estimate remains the cross-sector
product rule from Sec.~\ref{eq:cross_sector}:
\begin{equation}
  |V_{us}|_{\mathrm{hol}} = \sqrt{m_\mu/m_\tau} = 0.244\,,
\end{equation}
which is 8.7\% above the PDG value $0.2243$.
This arises not from meson mass matrix corrections but from
the holomorphic structure of the quark--lepton matching.

\paragraph{Cabibbo angle = Cartan angle (Koide 1981).}
Koide's original composite model~\cite{Koide:1983qe} predicts
\begin{equation}
  \tan\theta_C = \frac{\sqrt{3}\,(\sqrt{m_\mu} - \sqrt{m_e})}
  {2\sqrt{m_\tau} - \sqrt{m_\mu} - \sqrt{m_e}}\,,
  \label{eq:koide_cabibbo}
\end{equation}
which is \textbf{algebraically identical} to the Koide angle~$\delta$
(the arctan of this ratio defines~$\delta$).
Numerically $\theta_C = \delta = 0.22223$~rad, matching the PDG value
$|V_{us}| = 0.2243$ to 1.2\%.
The zero-parameter Wolfenstein hierarchy
$|V_{us}| = \delta$, $|V_{cb}| = \delta^2 = 0.049$, $|V_{ub}| = \delta^3$
gives $|V_{cb}|$ to 17\%.
The residual gap is absorbed by the structural identification
$A = 1/(1+\delta) = 9/11$, giving a zero-parameter Wolfenstein hierarchy:
\begin{equation}
  |V_{us}| = \frac{2}{9}\,,\qquad
  |V_{cb}| = \frac{4}{99}\,,\qquad
  A = \frac{9}{11}\,,
  \label{eq:wolfenstein_zero}
\end{equation}
matching PDG to $1.2\%$, $3.4\%$, and $1.0\%$ respectively.
The formula $A = 1/(1+\delta)$ implies an adjoint mass
$\mu_\Phi = m_\tau(1 + 1/\delta) = 11\,m_\tau/2 \approx 9.8\GeV$.

\paragraph{Origin of the Cartan angle.}
The angular potential generated by the $\mathbb{Z}_6$ discrete symmetry
of $\SU(3)_F$ with $\Nf = 3$ takes the form
\begin{equation}
  V(\delta) = -\cos(3\delta) + \alpha\,\cos(6\delta)\,.
  \label{eq:angular_potential}
\end{equation}
The minimum satisfies $\cos(3\delta) = 1/(4\alpha)$.
For $\alpha = 1/\pi$, the minimum occurs at
$\delta = 0.22249$~rad, which matches $2/9 = 0.22222$
to~$0.12\%$.
The coefficient $\alpha = 1/\pi$ is expected if the second
harmonic arises from the one-loop fluctuation determinant
around the instanton generating the first harmonic.
This provides a concrete, calculable mechanism for determining
the Cartan angle from the dynamics of $\SU(3)_F$ confinement.

\paragraph{The $1/N$ expansion.}
Setting $\delta = 2/N^2$ with $N = 3$,
every parameter in Table~\ref{tab:1overN} is a rational function of~$N$.
The lepton mass hierarchy, the CKM matrix, and the PMNS matrix
are all the \textbf{$1/N$ expansion of $\SU(3)_F$}.

\begin{table}[h]
\centering
\begin{tabular}{llll}
\hline
Parameter & Formula & Value & Order \\
\hline
$\delta$ (Cartan) & $2/N^2$ & $2/9$ & $O(1/N^2)$ \\
$K$ (Koide) & $j(j{+}1)/N$ & $2/3$ & $O(1/N)$ \\
$|V_{us}|$ & $2/N^2$ & $2/9$ & $O(1/N^2)$ \\
$A$ (Wolfenstein) & $N^2/(N^2{+}2)$ & $9/11$ & $1{-}O(1/N^2)$ \\
$|V_{cb}|$ & $4N^2/[N^2(N^2{+}2)]$ & $4/99$ & $O(1/N^4)$ \\
$\theta_{12} - \pi/4$ & $-2/N^2$ & $-2/9$ & $O(1/N^2)$ \\
$\theta_{23} - \pi/4$ & $+2/N^3$ & $+2/27$ & $O(1/N^3)$ \\
$\theta_{13}$ & $2j(j{+}1)/N^3$ & $4/27$ & $O(1/N^3)$ \\
\hline
\end{tabular}
\caption{All mixing parameters as the $1/N$ expansion of $\SU(N)_F$
  at $N = 3$, $j = (N{-}1)/2 = 1$.
  Scanning $N = 2,\ldots,10$: $\chi^2(\mathrm{PMNS}) = 0.5$ at $N=3$
  versus $604$ ($N{=}2$), $218$ ($N{=}4$). $N = 3$ is uniquely selected.}
\label{tab:1overN}
\end{table}

\paragraph{PMNS from the same angle.}
The Cartan angle also determines the PMNS deviations from tribimaximal
mixing:
\begin{equation}
  \theta_{12} = \frac{\pi}{4} - \delta + \frac{\delta^2}{N}\,,\qquad
  \theta_{23} = \frac{\pi}{4} + \frac{\delta}{3}\,,\qquad
  \theta_{13} = \frac{2\delta}{3} = \frac{4}{27}\;\mathrm{rad}\,,
  \label{eq:pmns_delta}
\end{equation}
giving $33.2^\circ$ ($0.7\%$ off), $49.2^\circ$ ($0.1\%$ off), and
$8.5^\circ$ ($0.9\%$ off) respectively.
The solar angle satisfies quark--lepton
complementarity~$\theta_{12} + \theta_C = \pi/4$; the reactor angle
is $\theta_{13} = K\cdot\delta$ where $K = 2/3$ is the Koide value.
Combined with Eq.~(\ref{eq:wolfenstein_zero}), six mixing parameters
follow from one number: $\delta = 2/9$.

\paragraph{Why CKM is hierarchical while PMNS is not.}
The leptons are \emph{exact} composites of $\SU(3)_F$:
their mixing (PMNS) receives unsuppressed non-perturbative
K\"ahler corrections $\sim O(1)$.
The quarks are \emph{partially} composite:
their mixing (CKM) is proportional to
$Y_i Y_j / (16\pi^2)$, hence suppressed.
This structural difference---not fine-tuning---explains
the observed hierarchy
$\theta_{\mathrm{PMNS}} \sim O(1)$,
$\theta_{\mathrm{CKM}} \sim \lambda \approx 0.22$.


%======================================================================
\section*{Acknowledgments}
%======================================================================

The author thanks M.~Porter for identifying Seiberg duality as
the mechanism for the sBootstrap during extended discussions
(2011--2024), C.~Brannen for the circulant matrix parametrisation
of the Koide formula, and the late M.~Sheppeard (Kea) for
early contributions connecting Koide to representation theory.
During the preparation of this work the author used Claude
(Anthropic) for algebraic verification, numerical checks,
and drafting assistance.
The author reviewed and edited all content and takes full
responsibility for the publication.


%======================================================================
\begin{thebibliography}{99}

\bibitem{Chew:1962}
  G.F.~Chew,
  Rev.\ Mod.\ Phys.\ \textbf{34}, 394 (1962).

\bibitem{Miyazawa:1966}
  H.~Miyazawa,
  Prog.\ Theor.\ Phys.\ \textbf{36}, 1266 (1966).

\bibitem{CattoGursey:1985}
  S.~Catto and F.~G\"ursey,
  Nuovo Cim.\ A \textbf{86}, 201 (1985).

\bibitem{Rivero:2024epjc}
  A.~Rivero,
  Eur.\ Phys.\ J.\ C \textbf{84}, 1058 (2024),
  arXiv:2407.05397.

\bibitem{VolkovAkulov:1973}
  D.V.~Volkov and V.P.~Akulov,
  Phys.\ Lett.\ B \textbf{46}, 109 (1973).

\bibitem{FradkinShenker:1979}
  E.~Fradkin and S.H.~Shenker,
  Phys.\ Rev.\ D \textbf{19}, 3682 (1979).

\bibitem{Kaplan:1991dc}
  D.B.~Kaplan,
  Nucl.\ Phys.\ B \textbf{365}, 259 (1991).

\bibitem{NelsonStrassler:2000}
  A.E.~Nelson and M.J.~Strassler,
  JHEP \textbf{09}, 030 (2000),
  arXiv:hep-ph/0006251.

\bibitem{Seiberg:1994bz}
  N.~Seiberg,
  Nucl.\ Phys.\ B \textbf{435}, 129 (1995),
  arXiv:hep-ph/9411149.

\bibitem{CsakiSchmaltzSkiba:1997}
  C.~Cs\'aki, M.~Schmaltz, and W.~Skiba,
  Phys.\ Rev.\ Lett.\ \textbf{78}, 799 (1997),
  arXiv:hep-th/9610139.

\bibitem{Tong:SUSY}
  D.~Tong,
  \textit{Supersymmetric Field Theory},
  lecture notes, University of Cambridge (2022),
  \url{http://www.damtp.cam.ac.uk/user/tong/susy.html}.

\bibitem{ISS:2006}
  K.~Intriligator, N.~Seiberg, and D.~Shih,
  JHEP \textbf{04}, 021 (2006),
  arXiv:hep-th/0602239.

\bibitem{SeibergWitten:1994}
  N.~Seiberg and E.~Witten,
  Nucl.\ Phys.\ B \textbf{426}, 19 (1994),
  arXiv:hep-th/9407087.

\bibitem{HMS:1997}
  T.~Hirayama, N.~Maekawa, and S.~Sugimoto,
  arXiv:hep-th/9705069 (1997).

\bibitem{Csaki:2011xn}
  C.~Csaki, Y.~Shirman, and J.~Terning,
  Phys.\ Rev.\ D \textbf{84}, 095011 (2011),
  arXiv:1106.3074.

\bibitem{Koide:1983qe}
  Y.~Koide,
  Phys.\ Lett.\ B \textbf{120}, 161 (1983).

\bibitem{Koide:1990}
  Y.~Koide,
  Mod.\ Phys.\ Lett.\ A \textbf{5}, 2319 (1990).

\bibitem{GellMannOakesRenner:1968}
  M.~Gell-Mann, R.J.~Oakes, and B.~Renner,
  Phys.\ Rev.\ \textbf{175}, 2195 (1968).

\bibitem{Sumino:2009}
  Y.~Sumino,
  JHEP \textbf{05}, 075 (2009),
  arXiv:0812.2103.

\bibitem{PDG:2024}
  S.~Navas \textit{et al.} (Particle Data Group),
  Phys.\ Rev.\ D \textbf{110}, 030001 (2024).

\bibitem{MartinRobertson:2019}
  S.P.~Martin and D.G.~Robertson,
  Phys.\ Rev.\ D \textbf{100}, 073004 (2019),
  arXiv:1907.02500.

\bibitem{Rivero:2005b}
  A.~Rivero,
  arXiv:hep-ph/0510068 (2005).

\bibitem{GattoSartoriTonin:1968}
  R.~Gatto, G.~Sartori, and M.~Tonin,
  Phys.\ Lett.\ B \textbf{28}, 128 (1968).

\bibitem{ORaifeartaigh:1975}
  L.~O'Raifeartaigh,
  Nucl.\ Phys.\ B \textbf{96}, 331 (1975).

\bibitem{Unsal:2008}
  M.~\"Unsal,
  Phys.\ Rev.\ Lett.\ \textbf{100}, 032005 (2008),
  arXiv:0708.1772.

\bibitem{Shih:2008}
  D.~Shih,
  JHEP \textbf{02}, 091 (2008),
  arXiv:hep-th/0703196.

\bibitem{HananyWitten:1997}
  A.~Hanany and E.~Witten,
  Nucl.\ Phys.\ B \textbf{492}, 152 (1997),
  arXiv:hep-th/9611230.

\bibitem{Nelson:1984}
  A.E.~Nelson,
  Phys.\ Lett.\ B \textbf{136}, 387 (1984).

\bibitem{Barr:1984}
  S.M.~Barr,
  Phys.\ Rev.\ Lett.\ \textbf{53}, 329 (1984).

\bibitem{Fayet:1976}
  P.~Fayet,
  Nucl.\ Phys.\ B \textbf{113}, 135 (1976).

\bibitem{MasieroVeneziano:1985}
  A.~Masiero and G.~Veneziano,
  Nucl.\ Phys.\ B \textbf{249}, 593 (1985).

\bibitem{LutyMohapatra:1997}
  M.A.~Luty and R.N.~Mohapatra,
  Phys.\ Lett.\ B \textbf{396}, 161 (1997),
  arXiv:hep-ph/9611343.

\bibitem{Kostant:1959}
  B.~Kostant,
  Amer.\ J.\ Math.\ \textbf{81}, 973 (1959).

\end{thebibliography}

\end{document}
